% 本文件由 LaTeX 排版 Agent 根据有效 translation.json 自动生成。
% author=Jerome; work_id=3004; title=致帕玛丘：驳耶路撒冷的若望

\chapter{致\Person{帕玛丘}：驳\Place{耶路撒冷}的\Person{若望}}

% structure: p0001 source=h1 target=omitted reason=page-title-already-rendered-by-parent

这篇驳斥\Place{耶路撒冷}的\Person{若望}的信函写于约398年或399年，是奥力振争议的产物。其直接起因是394年\Place{塞浦路斯岛}\Place{撒拉米}的主教\Person{埃皮法尼乌斯}访问\Place{耶路撒冷}。这位主教在\Place{复活堂}（§11）讲了一篇尖锐的反对奥力振主义的道理，这篇道理被认为直接针对\Person{若望}，以致后者派他的总执事去规劝讲道者（§14）。经历了诸多不体面的场面之后，\Person{埃皮法尼乌斯}建议\Person{热罗尼莫}及其朋友们与他们的主教分离（§39）。但他们如何能获得教会的服务呢？\Person{埃皮法尼乌斯}解决了这个困难：他把\Person{热罗尼莫}的兄弟带到他在\Place{埃琉特罗波利斯教区}\Place{阿德}所建的隐修院，在那里违背他的意愿祝圣了他，甚至使用武力克服他的反对（\Person{热罗尼莫}，\WorkTitle{书信}第五十一卷·1）。\Person{埃皮法尼乌斯}试图为他的行为辩护（\Person{热罗尼莫}，\WorkTitle{书信}第五十一卷·2），但\Person{若望}过了一段时间后向\Place{亚历山大城}上诉，指控\Person{热罗尼莫}及其支持者为分裂者。主教\Person{德奥斐洛}立即站在\Person{若望}一边：但他的使者\Person{依西多禄}写给\Person{若望}的一封信落到了\Person{热罗尼莫}手中（§37）。这封信表明\Person{依西多禄}只是作为\Person{若望}的党派人士而来，因此\Person{热罗尼莫}对信和送信人都暗含蔑视。争议就这样持续了大约四年，在几次和解尝试和表现出许多敌意（实际上等于将\Person{热罗尼莫}及其朋友们逐出教会）之后，争议停止了——也许是\Person{德奥斐洛}制止的，也许是受\Person{默拉尼雅}的影响。写给\Place{罗马}的\Person{帕玛丘}反对\Person{若望}的信（397或398年）突然中断，几乎可以肯定它在\Person{热罗尼莫}生前从未发表。后来\Person{热罗尼莫}对\Person{德奥斐洛}有如此大的影响，以至于我们发现他为\Person{若望}求情，\Person{若望}当时已经失宠于这位主教（\WorkTitle{书信}第八十六卷·1）。\par

本论文的日期存在争议。在§1中，\Person{热罗尼莫}说他写作于\Quote{三年之后}，即从\Person{埃皮法尼乌斯}访问\Place{耶路撒冷}（394年）算起，这应得397年。同样地，在§14中，他说\Person{埃皮法尼乌斯}对自己的冤屈思考了三年。另一个时间标记见于§43中的话，\Person{若望}\Quote{最近}试图从\Quote{那只威胁全世界脖颈的野兽}——即于395年底去世的普雷费克特·\Person{鲁菲努斯}——那里获得针对\Person{热罗尼莫}的流放判决。所有这些陈述都指向397年。另一方面，在§17中，他谈到他的\BibleBook{}{训道篇}和\BibleBook{}{厄弗所书}的\Quote{注释}是\Quote{\OriginalTerm{大约}{ferme} 十年前}写的；而\BibleBook{}{训道篇}的序言说，他曾在\Place{罗马}与\Person{布莱西拉}一起研读\BibleBook{}{训道篇}\Quote{\OriginalTerm{大约}{ferme} 五年前}，因此，是在写作本论文的十五年前。\Person{布莱西拉}死于384年。因此，研读\BibleBook{}{训道篇}可能发生在383年。而十五年将我们带到398年。另外，在§41中，\Person{热罗尼莫}对\Person{若望}说：\Quote{你似乎沉睡了十三年}，暗示\Person{若望}所抱怨的状况——即隐修士们出现在他的教区，或者至少他们离开了自己的教区——已经存在了十三年。\Person{热罗尼莫}于385年底或386年初离开了他被祝圣的教区\Place{安提约基雅}；因此，这十三年将我们带到399年，即\Person{瓦拉西}采用的日期。不过，在\WorkTitle{拉乌斯史}第117章中有一个暗示：\Person{鲁菲努斯}的朋友\Person{默拉尼雅}在\Quote{近四百名跟随\Person{保利诺}的隐修士的分裂}一事中提供了帮助，这被承认与\Place{伯利恒}的分裂有关，该分裂由\Person{保利尼亚诺}的祝圣问题引发。我们知道\Person{默拉尼雅}和\Person{鲁菲努斯}于397年初离开\Place{耶路撒冷}，并且在离开之前，\Person{热罗尼莫}和\Person{鲁菲努斯}已经和好。因此，似乎最有可能的是，这篇充满对\Person{若望}（\Person{鲁菲努斯}的同工）如此敌意并且包含对\Person{鲁菲努斯}本人不敬暗示（§11：\Quote{你的朋友们，像狗一样龇牙咧嘴，翘起鼻子}，这是\Person{热罗尼莫}对\Person{鲁菲努斯}一贯的描述）的论文，是在\Person{鲁菲努斯}与\Person{热罗尼莫}和好之前写的，即在386年底或387年初，并且因为它被中断并保持未发表，因为情况已经改变。\Person{瓦拉西}将其定在399年。他引用了支持较晚日期的段落，却奇怪地忽略了更明确支持较早日期的陈述。还应补充的是，\Person{热罗尼莫}写给\Person{德奥斐洛}的信（第八十二卷）显然是与本论文同时写的，并带有同样的情感；如果上述论证成立，那封信必须放在397年，而不是其前面注释所述的399年。那封短信（第八十六卷）写给\Person{德奥斐洛}的，在这种情况下很可能放在398年或399年，而不是那里所述的401年。\par

此篇论文是写给帕玛邱斯的信。帕玛邱斯因若望主教向罗马主教西里修控告热罗尼莫而感到不安。热罗尼莫开始（1）以必需攻击主教为由进行辩护。埃皮法尼乌斯指控主教为异端（2）。让主教明白地回答（3），因为只有骄傲（4）阻止他这样做。据说若望的解释或道歉信得到了德奥斐洛的认可（5）；但那封信并未触及关键问题，即对奥利振主义的控告。只处理了三点（6），而埃皮法尼乌斯提出了八点——即（7）奥利振的观点：（一）圣子不能看见圣父；（二）灵魂被禁锢在尘世的身体中，如同在监狱里；（三）魔鬼可能得救；（四）天主用以包裹亚当和厄娃的皮衣是人的身体；（五）复活后的身体没有性别；（六）对乐园的描述是寓言式的：树木代表天使，河流代表天上的美德；（七）穹苍之上和之下的水是天使和魔鬼；（八）天主的肖像在堕落中完全失落。若望没有针对第一点回答，而仅仅表达了他对天主圣三的信仰（8, 9），并且始终试图（10）将他与埃皮法尼乌斯之间的争论仅仅归结为保利尼安努斯被祝圣之事。热罗尼莫接着叙述了埃皮法尼乌斯与若望之间激烈争论的异常场面（11-14）。然后他转向奥利振主义的观念：天使被贬入人的灵魂（15, 16），人的精神进入天体（17），人的灵魂在先前就存在（18），并在受造物的等级中上下移动（19, 20）。若望没有在这些点上回答，而是满足于抗议摩尼教（21）。热罗尼莫就灵魂起源的问题（22）质问若望，鲁莽地倡导创造论。然后他转向复活后身体状态的问题（23），断言\Emph{肉身}将恢复如现在（24-27），基督（28）和我们都是如此，并从旧约（29-32）引用见证，并讨论主复活后的显现（34-36）。然后他转向详细审查若望写给德奥斐洛的信或\Quote{道歉书}（37），引用其措辞，并讲述伊西多尔受派遣的故事（37, 38），以及阿尔克拉乌斯伯爵试图讲和（39）。若望强调的保利尼安努斯被祝圣之事是一种托词（40, 41）。分裂是由于主教的有异端倾向，他到处被埃皮法尼乌斯谴责（42, 43）。\par

整封信语气激烈而轻蔑，傲慢地假设作者对他所写的困难主题拥有全部真理，并且他有权要求他的主教就他选择质问的每个点作信德宣认。其重要性在于：它在很大程度上决定了教会人士对于它所处理的点的信仰，以及处理所谓异端的方式，长达一千多年。\par

1. 如果按\BibleRef{罗 8:26}宗徒保禄所说，我们无法按着感觉祈祷，言语也无法表达我们自己内心的思想，那么判断别人的心，并追溯和解释他使用的特定词语和表达的含义，又该是多么危险呢？人的天性倾向于怜悯，在考虑别人的罪时，每个人都怜悯自己。因此，如果你责备一个在言语上犯罪的人，人会说他只是单纯；如果你指责一个人狡猾，你对他说话的人不会承认里面有更多的东西，只会说是不知，这样他就可以避免恶意的嫌疑。结果就会是，你，控告者，成了诽谤者，而被指责的一方不被视为异端者，而只是一个没有教养的人。你知道，帕玛邱斯，你知道，不是敌意或虚荣的欲望驱使我从事这项工作，而是我被你的信激励，并且出于我信德的热忱；如果可能，我希望所有人都明白，我不能被指责为急躁和鲁莽，因为我是在三年之后才发言。事实上，如果你没有告诉我许多人的心思因我将要讨论的\Quote{道歉书}而困扰，并疑惑不定，我本已决定继续保持沉默。\par

2. 所以，让诺瓦提安离开吧，他不肯向犯错之人伸出援手！让蒙丹和他的疯女人灭亡吧！蒙丹，他要把跌倒的人扔进深渊，使他们永不再起来。我们每天都犯罪，都会有些失误或别的什么。既然我们对自己仁慈，就不该对别人严苛；反而我们祈祷并恳求他，要么简单地指出我们自己的过错，要么公开地为别人的过错辩护。我不喜欢模棱两可；我不喜欢听到可以有双重含义的话。让我们以揭开的面纱\BibleRef{格后 3:18}瞻仰主的荣耀。从前以色列子民在两派之间摇摆不定。但厄里亚（意思是\Emph{主的勇者}）说：\BibleRef{列上 18:21}\Quote{你们摇摆不定，要到几时呢？如果主是天主，就跟随他；但如果巴耳是天主，就跟随他吧。}主自己也论及犹太人：\Quote{外邦的子孙对我说谎；外邦的子孙变得软弱，从他们的歧途上跛行。}如果实在没有理由怀疑他是异端（正如我所愿和所信的），那他为什么不用我自己的话来明确说出我的意见呢？他把这称为直率；我却视之为狡猾。他想让我相信他的信仰是纯正的；那么，他的话也应当是纯正的。的确，如果模棱两可只涉及一个词、一句话、或两三个地方，我可以因无知而宽容；我也不会用确定和清楚的标准来判断模糊和可疑之处。但事实上，这种\Quote{直率}不过是一种舞台伎俩，就像踮着脚尖在鸡蛋或谷穗上行走一样；到处都是怀疑和猜疑。你会以为他写的不是一篇信仰阐述，而是一篇关于某个虚构主题的辩论。他如今热衷的那一套，我们早就在学校里学过了。他穿上我们自己的盔甲来攻打我们。即使他的信仰正确，他说话也谨慎含蓄，但他过度的谨慎引起了我的怀疑。\BibleRef{箴 10:9}\Quote{行事正直的，步履稳健。}毫无缘由地背负恶名是愚蠢的。有人控告他一项他自己并不知情的罪名。让他自信地否认那悬于一词的指控，并自由地反过来指责他的对手。让他在反驳指控时，展现出与对方提出指控时同样的勇气。当他说完他所有想说的、并打算说的、以及无可怀疑的话之后，如果对手仍坚持诽谤，那就让他公开对簿公堂吧。我不希望任何人默默地忍受异端的指控，免得他若一言不发，他的不坦诚就会被那些不知其清白的人视为犯罪自觉的表现，尽管无需要求当事人到场，只要你手中握有他的信件，就能使他哑口无言。\newline{}\par

3. 我们都知道他写信给你说了什么，他指控了你什么，以及你认为他在其中诽谤了你什么。请逐一回答这些点；跟随这封信的足迹；不要忽略诽谤中的任何一个细节。因为如果你粗心大意，不小心遗漏了任何一点——我相信你曾发誓这样做——他会立刻叫喊：\Quote{现在，现在，你输了，一切都在这一点上。}同样的话在朋友和敌人耳中听来不一样。敌人会在灯芯草里找疙瘩；朋友却会把弯曲的也当作直的。世俗作家有句话说，情人在判断上是盲目的，不过也许你忙于研读圣经，无心留意这类文学。你绝不应夸耀朋友对你的看法。真正的见证来自敌人的口。相反，如果朋友为你说话，他会被视为审判官或同党，而非证人。你的敌人大概会这样说，也许他们不信你，只想惹恼你。但我是你声称从未故意伤害过的人，却总在你的信件中被你提来提去，我劝你要么公开宣告教会的信仰，要么按你所信的来说话。因为那种谨慎的斟酌措辞，无疑可以欺骗无知的人；但细心的听众和读者很快就会察觉陷阱，并会公开揭露那些推翻真理的地下坑道。亚略派（没有人比你更了解他们）很长一段时间假装他们谴责\Emph{同质}一词，是因为它引起反感，他们用蜜糖般的言语涂抹有毒的错误。但最后蛇展露了自身，它那隐藏在所有盘绕中的致命头颅，被圣神的利剑刺穿。教会，如你所知，欢迎悔改者，并被罪人的众多所压倒，以致为了迷途羊群的利益，不得不对牧者的创伤宽大处理。古代和现代的异端遵循同一规律——百姓听一套，司铎讲另一套。\newline{}\par

4. 首先，在我将这封信——您写给\Person{特奥斐罗}主教的信——译出并插入本书，并向您表明我理解您过分的谨慎和周详之前，我想先和您说几句规劝的话。您这种傲慢自大，使您拒绝回答那些就信仰向您提问的人，究竟是什么意思？您为何几乎把\Place{巴勒斯坦}那些拒绝与您共融的大批弟兄和隐修士团体视为公敌？天主子为了一只病羊，撇下了山上的九十九只，忍受了鞭打、十字架、鞭苔；祂背起负担，将那个沉溺于罪恶的淫荡女人耐心地扛在肩头，带到了天堂。难道您要扮演\Quote{至可敬的圣父}，挑剔的主教；凭您的财富和智慧、您的威严和学问，置身事外；傲慢地皱着眉头对待您的同仆，几乎不屑于看一眼那些被你们主的宝血所救赎的人？难道这就是您从宗徒的训诫所学的，要\BibleRef{伯前 3:15}\BibleQuote{时常准备好，对每一个问你们心中望德之理由的人，作出答复}？假如我们真的如您所说，寻找机会，并且以信仰热诚为借口，散布纷争，制造分裂，煽动争吵。那么，请将机会从那些盼望机会的人那里拿走；这样，您在信仰方面作出交代，解决了您所涉及的所有难题之后，就可以清楚地向所有人表明，争议不在于教义，而在于秩序。但也许当被问及信仰时，您出于深谋远虑而保持沉默，以免有人说您已证明自己是异端者——因为您向指控者作出了交代。若是如此，那么人们就不该反驳任何被指控的罪名，以免在否认之后被认定为有罪。我想，平信徒、执事和司铎们的指控不在您的眼里。因为您如同您不断夸耀的那样，可以在一个小时内制造一千名神职人员。但您必须答复我们的圣教父\Person{埃皮法尼乌斯}，他在他寄出的信中公开称您为异端者。当然，您在年纪、学问、圣德生活或全世界的判断方面都不如他。如果论年龄，您是一个年轻人给一位长者写信。如果论知识，您是一个不太有造诣的人给一位博学的人写信，尽管您的支持者坚称您比\Person{狄摩斯悌尼}更善辞令，比\Person{克里西普斯}更敏锐，比\Person{柏拉图}更智慧，或许他们也说服了您自己这是真的。至于他的生活和信仰热诚，我就不再多说了，以免显得我在寻求伤害您。当整个\Place{东方}（除了我们的圣教父\Person{亚大纳削}和\Person{保理诺}）被亚略异端和欧诺米异端侵扰时；当您不与\Place{西方}人共融时；在那最糟糕的流放中，这流放使他们成了宣信者，他，尽管只是一个简朴的修院司铎，却取得了\Person{欧提基乌斯}的信任，后来作为\Place{塞浦路斯}主教，\Person{瓦伦斯}没有骚扰他。因为他一直如此备受尊敬，以至于在位上的异端者认为，迫害这样的人只会给自己带来耻辱。因此，请您写信给他。回复他的信。让其余的人了解您的意图，并评判您的口才和智慧；不要将所有才华只留给自己。为什么当您在某一处受到挑战时，您却将武器转向另一处？在\Place{巴勒斯坦}向您提出了一个问题，您却在\Place{埃及}作答。当有些人眼睛模糊时，您却给没有受影响的人涂眼药。如果您告诉另一个人您本意是要给我们满足，那么这种行动完全出于骄傲；如果您告诉他我们没有要求的东西，那就完全多余。\par

5. 但您说：\Quote{亚历山大主教赞同了我的信。}他赞同了什么？您对\Person{亚略}、\Person{佛提诺}和\Person{摩尼}的正确言论。因为如今谁还指控您是亚略派？谁还把佛提诺和摩尼的罪过加在您身上？那些错误早已被纠正，那些敌人早已被粉碎。您还不至于愚蠢到公开维护一个您知道会冒犯整个教会的异端。您知道，如果您这样做了，您一定会立刻被免职，而您的心系于主教宝座的享乐。您如此调整您的措辞，以便既不惹恼愚钝之人，也不得罪您的支持者，他们无疑是标记着欺骗和油滑的；那么，对于剩下的五个，我们又当如何处理？因为它们没有提供模棱两可的机会。您写得很好，但文不对题。\Place{亚历山大}主教怎么知道您被指控的是什么，或者要求您宣认的是什么？您应该详细列出对您的指控，然后逐一应对。有一个古老的故事，说某人说话流畅，被滔滔不绝的言辞冲走，却不触及法庭上的问题，于是法官说出了睿智的评语：\Quote{妙！妙！但这所有妙处有何目的？}庸医对一切眼疾只有一种洗剂。一个人被指控许多事，而在化解指控时跳过一些，就等于承认了他所未提及的一切。您难道没有回复\Person{埃皮法尼乌斯}的信，并且自己选择了反驳的要点？无疑，在回复时，您依据的公理是，没有人勇敢到把剑架在自己脖子上。您喜欢哪种选择？您可以选择：您要么回复了\Person{埃皮法尼乌斯}的信，要么没有。如果您回复了，为什么对指向您的最重要、数量最多的指控置之不理？如果您没有回复，那么您那部在愚钝人中夸耀的\Quote{辩护书}，以及您大量散发给不了解此事之人的，又是什么呢？\par

6. 你要回答的问题，我很快会说明，共有八个方面。你只触及三个方面，然后跳过。至于其余方面，你保持意味深长的沉默。如果你坦诚地回答了七个，我仍然会坚持剩余的那项指控；而你未曾提及的，我会视为事实。但实际情况是，你抓住了狼的耳朵；既不能抓住不放，又不敢松手。你带着一种漫不经心的安全感和心不在焉的神情，掠过并触及了三个方面，其中没有什么或只有一点点重要内容。你的做法如此隐晦和严密，以至于你通过沉默所承认的，比你通过论证所反驳的还要多。每个人立刻就有权对你说，\BibleRef{玛 6:23}\BibleQuote{如果你内在的光成了黑暗，那黑暗是多么大啊！}即便在回答三个小问题时，你似乎说了些什么，但你并未摆脱嫌疑和责备；你的回答是模糊的，因此你无法欺骗听众，你宁愿保持完全的沉默，也不愿公开承认那些隐藏在阴暗中的东西吗？\par

7. 这些问题涉及\OriginalTerm{论原则}{\GreekText{Περὶ} \GreekText{Ἀρχῶν}}中的段落。第一个是这样的：\Quote{正如说子能看见父是不合适的，同样，认为圣神能看见子也是不合适的。}第二点是：灵魂被束缚在身体里如同牢狱；在人类被造于乐园之前，灵魂曾居住在天上理性受造物之中。因此，后来为了安慰自己，灵魂在\BibleBook{}{圣咏}中说：\BibleQuote{在我受磨难以前，我犯了错}；以及\BibleQuote{我的灵魂，回到你的安息吧}；\BibleQuote{引我的灵魂出离监牢}；以及在其他地方类似的话。第三，他说魔鬼和恶魔总有一天会悔改，并最终与圣人们一起掌权。第四，他将亚当和厄娃在堕落并被逐出乐园后所穿上的皮衣解释为人的身体，我们当然要假定此前在乐园里，他们没有肉体、没有筋、也没有骨头。第五，他极其明确地否认肉身的复活和身体的构造以及感官的区别，这既见于他对第一首\BibleBook{}{圣咏}的解释，也见于他的许多其他著作。第六，他以寓意解经的方式解释乐园，从而摧毁了历史真实性，将天使理解为树木，将天上的德能理解为河流，他通过寓意解释推翻了乐园历史中所包含的一切。第七，他认为在\WorkTitle{圣经}中所说的天上的水是神圣和超性的实体，而地上的和地下的水则是相反地属于魔鬼的实体。第八个问题是奥利振的诡辩，即人受造时所拥有的天主的肖像和模样已经丢失，在人被逐出乐园后就不在人身上了。\par

8. 这些就是刺伤你的箭矢；这些就是整封信中使你受伤的武器；或者我应该说，\Person{埃皮法尼乌斯}伏在你的膝前，将他的白发放在你的脚下，暂时放下主教的尊位，以这样的话祈求你的得救：\BibleQuote{求你赐给我和你自己以你得救的恩惠；拯救你自己，正如经上所载，脱离这扭曲的世代，\BibleRef{宗 2:40}并放弃\Person{奥力振}的异端以及一切异端，亲爱的。}往下又说：\Quote{在捍卫异端时，你激起对我的仇恨，并毁灭了我对你的爱；以致你甚至会使我们后悔曾与你共融，而你却如此坚决地捍卫\Person{奥力振}的错误与教义。}告诉我，辩论之王，你对八个段落中的哪一个作出了答复？目前，我对其他部分暂不置言。以第一个亵渎之词来说——即圣子不能见圣父，圣神不能见圣子。你的什么武器刺穿了它？我们得到的答复是：\Quote{我们相信神圣可钦崇的天主圣三同体，同永远，同光荣，同神性，并诅咒那些说在天主圣三的神性中有任何大小、不平等或任何可见之物的人。但正如我们说圣父是无形的、不可见的、永恒的，同样我们也说圣子和圣神是无形的、不可见的、永恒的。}你若不说这话，你就不属于教会。我不问是否曾有一段时间你拒绝说这话。我不讨论你是否喜爱宣讲这种教义的人；当那些人因表达这些观点而遭受流放时，你站在哪一边？或者，当司铎\Person{德奥}在教会中宣讲圣神是天主时，那个堵住耳朵、激动地冲出门外以免听到那亵渎之词的人是谁？我承认一个人，可以说，是信徒中的一员，即使他的悔改来得太晚。那个不幸的\Person{普雷泰克斯塔图斯}，在被选为执政官后去世，一个亵渎神的人和偶像崇拜者，惯于在戏谑中对真福\Person{达玛苏斯教宗}说：\Quote{立我为罗马主教，我立刻就成为一个基督徒。}你为什么费尽心思用许多言辞和复杂的句子来向我表明你不是亚略派？要么否认被指控的人说了所指控的话，要么如果他确实说出了这样的观点，就谴责他这样说。你还需要学习正统信仰者的热忱有多强烈。听宗徒的话：\BibleRef{迦 1:8}\BibleQuote{即使是我或从天上来的天使，给你们宣讲了与我们先前所宣讲的不同的福音，当受诅咒。}你试图减轻过失并隐藏有罪者的名字：仿佛一切正常，没有人被控亵渎，你以造作的语言构制了一份不必要的信仰宣认。你立即说出来，让你的信这样开头：\Quote{愿那胆敢写出这样东西的人受诅咒。}纯洁的信仰不容迟延。蝎子一出现，就必须被踩在脚下。达味，被证明是合天主心意的人，说：\BibleQuote{上主，我岂不恨那恨你的人？我岂不因你的敌人而憔悴？我以完全的恨恨了他们。}如果我听见我的父亲、母亲或兄弟对我的主基督说这样的话，我会像打疯狗一样打碎他们亵渎的嘴巴，我的手将是第一个举起来攻击他们的。他们对父亲和母亲说：\BibleRef{申 33:9}\BibleQuote{我们不认识你们，}这些人遵行了主的旨意。\BibleRef{玛 10:37}\BibleQuote{爱父亲或母亲超过基督的人，不配属于祂。}\par

9. 据称你的老师——你称他为天主教徒，并坚决捍卫他——说过：\Quote{圣子不见圣父，圣神不见圣子。}你告诉我圣父是不可见的，圣子是不可见的，圣神也是不可见的，仿佛天使，包括革鲁宾和色辣芬，按照其本性，对我们的眼睛也是不可见的。达味甚至对天的景象也存疑：他说：\BibleQuote{我要看见你的手指所造的天。}我要看见，不是我看见。我要在揭开帕子时看见主的荣耀；但现在\BibleRef{格前 13:9}\BibleQuote{我们只知道一部分，只预言一部分。}问题是圣子是否看见圣父，而你却说：\Quote{圣父是不可见的。}争论的是圣神是否看见圣子，而你回答：\Quote{圣子是不可见的。}争议点是天主圣三是否彼此互见；人的耳朵无法忍受这样的亵渎，而你却说天主圣三不可见。你在其他一切方面游走于赞美之域；你在无人想听的事情上浪费口才。你转移听众的注意力，以避免告诉我们所问的事。但就算所有这些是多余的，我们承认你不是亚略派，甚至从未是过。我们容许在解释第一个段落时，你毫无嫌疑，你所说的一切都是坦诚的，没有错误。我们同样坦诚地对你说：我们的天主之人\Person{埃皮法尼乌斯}曾指控你是亚略派吗？他曾把\Emph{无神的}\Person{欧诺弥}或\Person{阿厄里}的异端归咎于你吗？整封信的要点是你追随\Person{奥力振}的错误教义，并与他人同流合污于这异端。为什么当你被问及一点时，你却回答另一点；好像在对愚人说话，隐瞒信中的指控，并告诉我们你在\Person{埃皮法尼乌斯}面前在教会里说的话？有人要求你作信仰宣认，你却用你非常雄辩的论述来折磨我们。我恳请我的读者记住主的审判台，既然你们知道你们将按你们的审判受审，就不要偏袒我或我的对手，不要看辩论者的身份，而要看事实本身。让我们继续我们开始的事。\par

10. 你在信中写道，在\Person{保利尼亚诺}被祝圣为司铎之前，\Person{厄丕法尼}主教从未因\Person{奥力振}的错误而责备你。首先，这一点是可疑的，我必须考虑两人中该相信谁。他说他曾反对，你否认；他提出证人，你拒绝听取他们作证；他甚至说除你之外还有另一人被控诉：你两件事都不承认；他通过一位圣职人员寄信给你，并要求答复：你保持沉默，不敢开口，在巴勒斯坦被挑战却到亚历山大发言。该相信谁，不是我能说的。我想你自己也不敢在如此卓越的人面前，坚持自己为真而指责他为假。但可能各自从自己的角度说话。我要叫一个证人反对你，那证人就是你自己。因为如果教义上没有争议，如果你没有激怒一位老人的怒气，如果他未曾给你答复，你（并不擅长言辞）何必在一次讲道中向大众，也向这样一位卓越的人讨论全部教义——天主圣三、主的身体降孕、十字架、地狱、天使的本性、灵魂的状况、救主的复活和我们自己的复活，并且是在此世上（你的手稿或许遗漏了这些主题）？并且如此自信地一口气奔驰而过，不停歇？我们该如何说教会的古代作家，他们用许多卷册也难以解释单个难题？我们该如何说那蒙选之器、福音的号角、我们狮子的咆哮、外邦人的雷霆、基督教雄辩之河，当他面对\BibleRef{哥 1:26}那历代以来隐藏的奥秘，以及\BibleRef{罗 11:33}天主智慧和知识的渊深时，更多的是惊叹而非讨论？我们该如何说那预指童贞女的\Person{依撒意亚}？那一件事已使他力不能及，他说：\BibleRef{依 53:8} \Quote{谁能述说他的世代？}在我们的时代，竟出现了一个可怜的小人，他用舌头一转，以超过太阳的光辉，谈论所有教会问题。如果没有人要求你展示，一切平静，你自愿进入如此危险的讨论实属愚蠢。另一方面，如果你讲话的目的是为了满足你对信仰的责任，那么冲突的原因就不是某位司铎的祝圣（显然那是很久以后的事）。你只欺骗了不在场的人，你的信只取悦了远方人的耳朵。\par

11. 我们当时在场（我们知道全部情况），当\Person{厄丕法尼}主教在你的教会中反对\Person{奥力振}时，他表面上是攻击他，实际上是你。你和你的同伙像狗一样咧嘴，缩起鼻子，挠头，互相点头，说那\Quote{愚蠢的老头子}。你难道不是在主墓前派你的总执事叫他停止讨论这些事吗？有哪位主教曾在民众面前向自己的司铎下过这样的命令？当你从复活堂前往圣十架堂时，各年龄、男女的人群涌向他，把自己的小孩献给他，亲吻他的脚，扯他的衣边，他寸步难行，几乎无法在人潮的波涛中站立，你难道不是被嫉妒折磨，以致喊出\Quote{那虚荣的老头子}吗？而且你竟不羞耻当面对他说，他停下来是故意为之。请回想那一天，被召集的人民仅仅因为希望听到\Person{厄丕法尼}，就一直等到第七时辰，以及你当时演讲的主题。你果然愤怒地抨击\OriginalTerm{拟人论者}{Anthropomorphites}，他们以淳朴的幼稚认为天主真的有圣经中读到的肢体；你用眼睛、双手和每一个手势表明你指的是那老人，并希望他被怀疑犯有那最愚蠢的异端。当你筋疲力尽，口舌干燥，头往后仰，嘴唇颤抖，在全体渴望结束的人民的满意中终于结束时，那疯狂愚蠢的老头子如何对待你？他站起来表示要说几句话，先以声音和手向会众致意，然后说道：\Quote{凡是由一位在主教职上是我兄弟、但在年龄上是我儿子的人，针对拟人论者异端所说的一切，都说得善而忠信，我的声音也谴责那异端。但既然我们谴责这异端，也应同样谴责奥力振的乖谬道理。}我想你不会忘记随后爆发出的笑声和欢呼声。这就是你在他信中所谓的\Quote{他随意向人民说任何话，不论是什么。}他果然疯了，因为他在你的地盘上反驳了你。\Quote{随意说任何话，不论是什么。}要么你称赞他，要么责备他。为什么你在这里和其他地方都如此摇摆不定？如果他说的好，为何不公开宣扬？如果坏，为何不大胆谴责？然而，让我们注意，这位真理与信仰的支柱，竟敢说如此杰出的人向人民随意说话，他自己以何等的智慧、谦虚和谦逊言及自身：\Quote{有一天我在他面前说话；借着当日读经中的一些话，我在他及全体教会面前，表达了关于信仰及教会全部教义的观点，这些观点我靠天主的恩宠在教会内和我的要理讲授中不断教导。}\par

12. 请问，这种厚颜无耻和夸夸其谈是什么意思？所有哲学家和演说家都攻击莱昂提尼的高尔吉亚，因为他竟敢公开承诺回答任何人可能向他提出的任何问题。如果司铎的荣誉和对您头衔的尊重没有约束我，如果我不知道宗徒所说的：\Quote{弟兄们，我不晓得他是大司祭；因为经上记载：不可毁谤你百姓的首领。}，那么我将会如何愤怒地抱怨您所讲述的事！相反，您却在言语和行为上轻视几乎是整个主教职之父、昔日圣德纪念碑的人物，从而贬低了您头衔的尊严。您说，在某一天，当日课中的某件事激起了您，您当着他的面，并在整个教会面前，就信仰和教会的一切教义发表了一篇演讲。此后，我们不得不惊叹德摩斯梯尼的薄弱；因为我们得知，他花了很长时间精心构思他那篇反对埃斯基涅斯的精彩演说。我们仰望图留斯是完全错误的；因为根据在场的科尔内利乌斯·内波斯的说法，他的功绩仅仅在于，他几乎逐字逐句地发表了著名的为煽动护民官科尔内利乌斯辩护的演说。看哪，一个吕西亚斯和一个格拉古为我们兴起了！或者，举一个更近代的人，昆图斯·阿特留斯，他所有的能力都像一笔现成的钱，需要有人告诉他何时停止，凯撒奥古斯都都对他说得很好：\Quote{我们的朋友昆图斯必须刹住车。}\par

13. 有哪个神志清醒的人会宣称在一篇讲道中讨论了信仰和教会的一切教义？请告诉我，哪一课经文如此充满圣经的全部韵味，以至于它在礼仪中的出现促使您进入竞技场，将您的才智置于险境？如果您没有被您滔滔不绝的雄辩所淹没，您或许会相信，未经深思熟虑而谈论全部教义是不可能的。但事实如何？您承诺一件事，却呈现另一件。我们的习惯是，在四十天期间，向那些将要受洗的人公开讲授关于圣而可敬的天主圣三的教义。如果当天的课经刺激您在一小时内讨论所有教义，那么重复之前四十天的教导又有何必要？但如果您意在总结整个四旬期所说的话，那么某一天的一节课经如何能\Quote{刺激}您谈论所有这些教义？然而即使在这里，他的语言也是模棱两可的；因为他可能从那一节特定的课经出发，简要地回顾了他通常在四旬期四十天里在教会中向慕道者所讲授的内容。因为无论是用许多话说很少的事，还是用很少的话说很多的事，同样都是雄辩。还有另一种可允许的含义：一旦那一节课经给了他刺激，他就被如此热烈的演说激情所点燃，以至于四十天来从未停止说话。但是，那位悠闲的老人，一直专心听他说话，渴望知道从未听过的事，也几乎要从座位上睡着了。然而，我们必须忍受；也许，这又是一个我们熟悉的他那种质朴的例子。\par

14. 让我们引用其余部分，在那里，经过他那令人困惑的讨论迷宫之后，他表达得毫不含糊，而是公开地，并以这种方式结束了他奇妙的讲道：\Quote{当我们当着他的面如此说后，出于对他极度的尊敬，我们邀请他随后发言，他赞扬了我们的宣讲，说他对此感到惊奇，并向所有人宣布这是公教信仰。}您对他极度的尊敬体现在您对他极度的侮辱中，当时您通过总执事命他保持沉默，并大声宣称是爱慕虚荣使他在人群中逗留。现在的事是过去的钥匙。从那以后的整整三年里，他默默地沉思他所遭受的冤屈，摒弃一切个人争执，只请求您更准确地表达您的信仰。而您，有无尽的资源，利用全世界的宗教牟利，一直在派遣那些非常高贵的使节到处奔走，并试图唤醒那位老人，让他回答您。事实上，既然您给了他如此显著的荣誉，他赞扬您的言论，尤其是那些即席发表的言论，是合宜的。但是，正如人们有时会赞扬他们并不赞同的东西，通过无意义的赞美来滋养他人的愚蠢一样，他不仅赞扬了您的言论，而且赞扬并惊奇于它们；更重要的是，为了夸大这种惊奇，他向全体人民宣布它们符合公教信仰。他是否真的说了所有这些，我们自己是见证人。事实是，他到我们这里来时，被您的话吓得半死，说他太贸然与您共融了。此外，当整个隐修院恳求他从白冷回到您那里，他无法抗拒这么多人的请求时，他确实在傍晚回来了，但只是为了在午夜再次逃走。他写给教宗西利修斯的信也证明了同样的事，如果您读了这些信，就会清楚地明白他是在什么意义上惊奇于您的言论并承认它们是公教的。但我们是在打谷糠，用许多话驳斥毫无根据的废话和老妇的寓言。\par

15. 让我们转到第二个要点。他在这里好像无需考虑什么，便漫不经心地空谈自夸，假装睡着，好让读者也昏昏欲睡。但我们正在谈论其他与信仰有关的事，即：所有可见与不可见的事物、天上的能力和地上的受造物，都由同一个创造者，即天主，也就是天主圣三所创造，正如蒙福的达味所说：\Quote{因天主的话，诸天确立；因祂口中的气，万军生成}；人的受造同样是明证，因为是天主亲自从地上取来泥土，借着祂自己嘘气的恩宠赐给它一个理性的、赋有自由意志的灵魂；这灵魂不是祂本性的一部分（如某些不敬者所教导的），而是祂自己的化工。关于圣天使，基督徒的信仰同样遵循圣经，圣经论到天主说：\Quote{祂使祂的天使成为风，使祂的仆役成为火焰。}圣经不允许我们相信他们的本性是不可改变的，因为经上说：\BibleRef{犹 1:6}\Quote{又有不守本位、离开自己住处的天使，主用永远的锁链把他们拘留在黑暗里，等候大日的审判}；因此，我们知道他们改变了，失去了自己的尊威和光荣，变得更像魔鬼。然而，关于人的灵魂是由天使的堕落或悔改而来的说法，我们从未相信，也从未如此教导（愿天主不许！），我们承认这种见解与教会的训导相悖。\par

16. 我们想知道，在人在乐园中受造以前、亚当由尘土形成以前，灵魂是否已在理性受造物中；它们是否拥有自己的等级、生命、延续和存在；以及奥利金的学说是否正确，他声称所有理性、无形、不可见的受造物，若变得懈怠，便逐渐沉沦到更低层次，并根据所降之处而取用身体（例如，起初可能是以太体的，后来是气态的）；当它们接近大地时，便穿上最粗重的身体，最后被束缚于人的肉体；而那些自甘堕落、与它们的首领魔鬼一同离弃天主事奉的魔鬼本身，若稍有悔改，便穿上人的肉体，以便在复活之后经历补赎，并经过与它们降到肉体时相同的回路，最终脱离气态和以太体的身体，返回与天主亲近；那时，一切在天上、地上和地下的膝都要向天主跪拜，使天主成为万物中的万有。这些才是真正的问题，你为何避而不谈，离开战场，把自己钉在遥远且毫不相关的讨论区域中？\par

17. 你相信是一个天主创造了所有可见与不可见的受造物。亚略若说万物是藉着圣子受造的，他也会承认这一点。若有人指控你信奉马尔基翁的异端，即引入两个天主，一个良善之天主，一个公义之天主，并声称前者创造不可见之物，后者创造可见之物，那么你的回答正好适合于使我在这种问题上满意。你相信创造宇宙的是天主圣三。亚略派和半亚略派否认这一点，他们亵渎地主张圣神不是创造者，而是受造者。但现在谁指控你是亚略派呢？你说人的灵魂不是天主本性的一部分，仿佛现在埃皮法纽指控你是摩尼教徒似的。你反对那些断言灵魂出自天使、认为它们堕落后的本性变成人类实体的人。不要隐瞒你所知道的，也不要装出你不具有的单纯。奥利金从未说过灵魂出自天使，因为他教导\Emph{天使}一词描述的是职务而非本性。在他的著作\OriginalTerm{论原理}{\GreekText{Περὶ} \GreekText{᾿Αρχῶν}}中，他说天使、上座者、宰制者、掌权者、管辖这世界和黑暗的，以及\BibleRef{弗 1:21}每一个有名可称的，不仅在今世，也在来世，都成为它们所穿上之身体的灵魂，这些身体是它们因自己的欲望或为所委任的职务而取用的；太阳本身、月亮以及所有星辰，都是曾经有理性和无形体的受造物的灵魂；它们虽暂时屈伏于虚幻之下，即烈火般的身体，我们因无知和缺乏经验而称之为世界的光体，但它们终将从败坏的奴役中得获释放，进入天主子女光荣的自由。因此，一切受造物一同叹息，同受产痛。宗徒也哀叹说：\BibleRef{罗 7:24}\Quote{我真苦啊！谁能救我脱离这死亡的身体？}现在不是争论这学说的时候，这学说是部分外教、部分柏拉图式的。大约十年前，我在\BibleBook{}{训道篇}的\Quote{注释}和\BibleBook{厄弗所}{书}的注释中，我想自己已经向有思想的人阐明了自己的观点。\par

18. 如今我恳求你——你的口才如此丰沛，在一次讲道中就能阐释关于所有议题的真理——请你用简洁明了的言辞回答你的质问者。当天主用泥土造人，并借着祂自己嘘气的恩宠赐予他灵魂时，那后来由天主的嘘气所赋予的灵魂，是否先前已经存在并且存续？它在哪里？还是它在第六日，当身体由泥土形成时，从天主的能力中获得了存在和生命的能力？你对此缄默不语，假装不知道别人想问什么，却忙于无关的问题。你对\Person{奥力振}避而不谈，却对\Person{马尔强}、\Person{阿波利纳里}、\Person{欧诺弥}、\Person{摩尼}以及其它异端者的荒谬大发雷霆。别人向你要手，你却伸出脚来，同时暗中暗示你所持守的教义。你对像我们这样纯朴的人说悦耳的话，但这样做丝毫不让你的同党不悦。\par

19. 你说魔鬼是由天使变成的，而非灵魂，好像你不知道按\Person{奥力振}的说法，魔鬼本身是属于气态身体的灵魂，若悔改，在成为魔鬼之后注定要变成人的灵魂。你写道天使是可变的；在虔诚意见的掩饰下，引入一种不虔，声称经过许多世代之后，灵魂不是从天使产生的，而是从天使最初变成的某种东西产生的。我想把我的意思说得更清楚：假设一个具有护民官等级的人因自己的行为不当而被降级，经过骑兵的多个级别直至成为列兵，他是否立刻不再是护民官而成为新兵？不，他先是上校，然后依次是二百夫长、百夫长、军需官、巡逻兵、骑兵，最后是新兵；尽管我们的护民官最终成为普通士兵，但他并非从护民官等级直接变成新兵，而是变成上校。\Person{奥力振}用\OriginalTerm{雅各伯天梯}{Jacob's ladder}来教导说，理性受造物逐渐下降到最低一级，即血肉之躯；并且不可能有人从第一百突然跌到第一而不经过中间的数字到达最后，如同从梯子的阶梯下降；他们在从天堂到地球的旅程中，每当更换停留地，就更换身体。这些就是你的伎俩和诡计，你藉此把我们说成是\Quote{佩路西人}和\Quote{负重的牲畜}和\Quote{属血气的人}，他们\Quote{不领受属神的事}。\BibleRef{格前 2:14} 你们是\Quote{耶路撒冷人}，甚至能嘲弄天使。但你们的奥秘正被拖到光天化日之下，你们的教义不过是异教寓言的杂烩，正在基督徒的耳中公开暴露。你们如此钦佩的东西，我们很久以前在\Person{柏拉图}那里发现时就已鄙视。我们鄙视它，因为我们接受了基督的愚妄。我们接受了基督的愚妄，因为\BibleRef{格前 1:25} 天主的愚笨总比人智明。对我们这些基督徒和天主的司铎来说，难道不应当感到羞愧，把我们的言辞纠缠在模棱两可的意义中，好像只是在开玩笑？保持词语在两个意思之间摇摆，欺骗说话者本人甚于听者？\par

20. 你们中的一人，当我逼问他对灵魂的看法——灵魂是否在肉体之前就已存在？——他回答说灵魂和身体是一同存在的。我知道这人是异端者，想要在言语上陷害我。最后我捉住他说，灵魂从它开始使身体有生命之时才获得此名，而它先前被称为魔鬼、撒殚的天使、淫乱的神，或者反过来，称为宰制、权力、神的代理人、使者。然而，若灵魂在\Person{亚当}被安置于乐园之前就已存在（无论何种等级和状态），并且生活和行动（因为我们不能认为无形而永恒的东西像睡鼠一样迟钝和麻木），那么必然有某种预先的原因，来解释最初没有身体的灵魂后来被赋予身体。如果灵魂没有身体是自然的，那么它在身体内就是违反自然的。如果在身体内是违反自然的，那么身体的复活就是违反自然的。但复活不会是违反自然的；因此，按你所说，那违反自然的身体在复活时将会没有灵魂。\par

21. 你说灵魂不是天主的本质。嘿！这正是我们所预料的，因为你谴责邪恶的摩尼，提起他的名字就是玷污。你说天使不会变成灵魂。我在某种程度上同意，尽管我知道你赋予这些词的意义。但是，既然我们已经知道了你所否认的，我们希望知道你所相信的。你说：\Quote{天主拿泥土塑造了人，并藉着祂自己嘘气的恩宠赐予他一个理性的灵魂，并藉着自由意志的恩宠，而非祂自己神圣本性的一部分（如某些人邪恶地主张的那样），而是祂自己的手工。}看他如何不厌其烦地对我们没有问及的事情夸夸其谈。我们知道天主用泥土造了人；我们知道祂向他的脸上嘘气，人就成了一个有生命的灵魂；我们并不无知，灵魂具有理性和自由选择，我们知道它是天主的工艺。没有人怀疑摩尼在说灵魂是天主的本质时是错误的。现在我问：那个灵魂，天主的工艺，以自由意志和理性为特征，并且不是造物主的本质，是什么时候被造的？它是在人由泥土造成、生命的气息吹入他脸上的同时被造的吗？或者，它先前已经存在，并与理性的、无形体的受造物相结合并生活，后来才被赋予天主的嘘气？对此你保持沉默；在这里你装出乡间的单纯，并用圣经的话来掩饰非圣经的主张。在你肯定没有人想知道的事情（即灵魂不是天主自己本性的一部分，如某些人邪恶地主张的那样）的地方，你本应宣布（这是我们所有人都想知道的）它并不是先前存在、祂先前创造、长久居住在理性的、无形体的、不可见的受造物中的东西。你一件也不说；你提出\Person{摩尼}，却把\Person{奥力振}隐藏起来，就像孩子们要吃东西时保姆用一些玩笑打发他们一样，你将我们这些可怜乡民的思绪引向别处，好让我们专注于舞台上的新角色，而不去要求我们想要的东西。\par

22. 但假设事实是你仅仅省略了这一点，而你的单纯并不意味着你足够精明而隐藏的东西。一旦你开始谈论灵魂，并从人的最初受造中推导出这个如此重要主题的论据，为什么你半途而废地放下讨论，突然转向天使和主身体存在的条件？为什么你要越过如此巨大的困难泥沼，让我们陷入泥淖？如果天主的嘘气（你并不喜欢的观点，且你现在悬而未决的一点）是创造人类灵魂，那么\Person{厄娃}从哪里得到她的灵魂，既然天主没有向她的脸上嘘气？但我不想在\Person{厄娃}身上停留，因为她作为教会的预像，是从她丈夫的一根肋骨造成的，在这么多世代之后，不应再受到她后代的诽谤。我问\Person{加音}和\Person{亚伯尔}，我们原祖父母的长子，他们的灵魂从哪里来？整个人类后代，我们应当怎样想他们灵魂的起源？它们是通过繁殖而来，如同野兽一样？所以，正如身体来自身体，灵魂也来自灵魂？还是说，理性的受造物渴望身体存在，逐渐降到地上，最后甚至被束缚在人的身体上？的确（正如教会依据救主的话\BibleRef{若 5:17}\BibleQuote{我父至今工作，我也工作}；以及\BibleBook{依撒意亚}{}的经文\BibleQuote{谁在他里面造了人的精神}；以及\BibleBook{圣咏}{}\BibleQuote{祂一一造了他们的心}所教导的）天主每天都在创造灵魂——祂，愿意就能做到，并且从不停止作为创造者。我知道你惯常对此说些什么，以及你如何用奸淫和乱伦来对抗我们。但关于这些的争论是冗长的，会超过我们有限的时间。同样的论点可以反驳你，并且任何在创造者当前安排中似乎不配的东西，再次并不不配，因为它是祂的恩赐。从奸淫中出生不会归咎于孩子，而是归咎于父亲。如同种子的情况，滋养的大地没有犯罪，被撒入犁沟的种子也没有犯罪，以及在谷粒发芽所影响下的热力和湿气也没有犯罪，而是某个人，比如窃贼和强盗，他用欺诈和暴力拔起种子。同样，在人的生育中，相当于大地的子宫接收自己的东西，滋养它所接收的，然后给它所滋养的东西一个身体，并把它所形成的身体分为几个肢体。在这些肚腹的隐秘处，天主的手一直在做工，那里有身体和灵魂同一个创造者。不要轻视你创造者的美善，祂按祂的意愿塑造了你，造了你。祂本身就是天主的德能和天主的智慧，在童贞女的子宫里为自己建造了房屋。被宗徒列入圣人中的\Person{依弗大}是一个妓女的儿子。但是听着：\Person{辣黑耳}生的\Person{厄撒乌}有\Person{依撒格}，一个\Quote{多毛的人}，在心思和身体上，像好的麦子，退化成毒麦和野燕麦；因为罪恶和美德的原因不在于种子，而在于被生者的意志。如果生来就具有人的身体是一种过错，那么\Person{依撒格}、\Person{三松}、\Person{洗者若翰}如何是恩许之子？我相信，你看到了坚持自己信念的勇气是什么。假设我错了，我公开说出我的想法。那么，你也同样要么自由地宣示我们的意见，要么坚定地持守你自己的。不要把你放在我的战线上，这样你通过假装单纯就安全了，并且你可以在你愿意的时候从背后刺伤你的对手。我现在不可能写一本书来反对\Person{奥力振}的意见。如果基督赐给我们生命，我们将另写一部著作来讨论它们。现在的问题是，被告是否回答了对他的提问，以及他的回答是否清晰明确。\par

23. 让我们从这一点转向最为人熟知的问题，即有关肉身与身体复活的问题；在此，我的读者，我要提醒你，要知道我是怀着对天主的敬畏和审判而说话，你也应当这样聆听。因为，如果在他的阐述中能找到纯正的信德，并且没有不忠的嫌疑，那么我不会愚蠢到去寻找指责他的机会，而当我想要因他人的过失而责备他时，自己却因诽谤而受责备。因此，我请你阅读以下关于肉身复活的内容；阅读之后，如果它令你满意（我知道它很能取悦无知者），请暂且不下判断，稍等片刻，暂不发表意见，直到我完成答复；如果在那之后它仍令你满意，那么你便可将诽谤的烙印定在我们身上。\Quote{我们也宣认祂在十字架上所受的苦难、祂的死亡与埋葬（这是世界的救恩），以及祂真实而非想象的复活；并且祂作为死者的首生者\BibleRef{哥 1:18}，将我们肉身的初果（在被置于坟墓后，祂使之复活）带到了天堂，从而藉着祂自己身体的复活给予我们复活的希望；因此，我们都希望像祂复活那样从死者中复活；并非藉着任何外在陌生的身体（那只是暂时假借的幻影），而是如同祂自己复活时所藉的那个被放置在咱们门口的圣墓中的身体一样，我们也要藉着我们现在所穿戴、现在所埋葬的同一身体，因着同一理由和同一命令而希望复活。因为就如宗徒所说\BibleRef{格前 15:44}：\InnerQuote{播种的是属生灵的身体，复活的是属神的身体}；关于他们，救主在教导中说\BibleRef{路 20:35}：\InnerQuote{那些堪得来世及从死者中复活的人，不娶也不嫁，因为他们不能再死，而且如同天使一样，因为他们是复活之子}。}\par

24. 再者，在他信函的另一部分，即在他自己讲道的结尾，为了欺骗无知者的耳朵，他关于复活大做文章、大声喧哗，但用的是模棱两可、平衡的语言。他说：\Quote{我们并未省略我们的主耶稣基督的第二次光荣来临，祂要在自己的光荣中来临，审判活人死人；因为祂要唤醒所有死人，使他们站立在祂的审判座前；并将按照每个人在肉身内所行的，无论是善是恶，予以报应；因为每个人都要因着在肉身内活出纯洁公义的生活而得到荣冠，或因着同时成为享乐与邪恶的奴隶而被定罪。} 我们在福音中读到，在世界的终结，如果可能，连选民也要被迷惑\BibleRef{玛 24:24}，我们看到在这段话中得到了应验。无知的群众听到死人、埋葬，听到真实而非想象的死人复活，听到我们肉身的初果在主的身内到达了天上，听到我们将要复活，不是在外来的、陌生的身体（那只是幻影），而是如同主在放置在咱们中间的圣墓中的身体一样复活，我们也将以现在所穿戴、所埋葬的同一身体在审判之日复活。而且，唯恐有人以为这还不够，他在最后部分补充说：\Quote{祂将按照每个人在肉身内所行的，无论是善是恶，予以报应：因为每个人都要因着他的纯洁公义生活而在身体内得到荣冠，或因着成为享乐与邪恶的奴隶而被定罪。} 听到这些，无知的群众不怀疑任何诡计，不怀疑在这关于死人、埋葬身体和复活的所有喧嚷中有什么陷阱。他们相信事情正如所说的那样。因为百姓耳中的虔诚比司铎心中的更多。\par

25. 读者啊，我一再劝你要忍耐，并学习我通过忍耐所学到的事；然而，在我揭开那龙的面纱，并简要解释奥利金关于复活的观点之前（因为除非你清楚看到毒药是什么，否则你无法知道解药的功效），我请你谨慎阅读他的陈述，并反复思考。要特别留意：虽然他九次提到身体的复活，却从未一次提及肉身的复活，你可以合理地怀疑他是故意省略的。奥利金在若干地方，尤其是在他的第四卷\WorkTitle{论复活}、\WorkTitle{第一篇圣咏的解释}和\WorkTitle{杂集}中，说在教会中有一种双重的错误，我们和异端者都卷入其中：\Quote{我们因着淳朴和对肉身的偏爱，说同样的骨头、血、肉，简言之，四肢、面貌和整个身体结构，在末日都要复活：以致我们果真要用脚行走，用手工作，用眼观看，用耳听闻，还带着一个永不满足的肚腹和一个消化食物的胃。因此，相信这一点，我们就说必须吃喝、履行自然的功能、娶妻生子。因为如果没有婚姻，生殖器官有什么用？如果食物不需要咀嚼，牙齿有什么用处？肚腹和食物又有什么好处，如果按照宗徒所说的，它们都要被消灭？同一宗徒又大声宣布：\BibleRef{格前 15:58} \BibleQuote{血肉不能承受天主的国，可朽坏的也不能承受不可朽坏的。}}据他说，这就是我们乡野淳朴之人所持守的。至于那些异端者，其中有玛尔吉安、阿佩肋、瓦伦提努、玛乃斯（玛尼亚的同义词），他说他们完全否认肉身和身体的复活，只允许灵魂得救，并认为我们说我们要按主的样式复活是徒劳的，因为我们的主自己也是以幻影的身体复活的，不仅他的复活，就连他的诞生也是\Emph{幻象说的}或想象出来的，也就是说，更多是表面的而非真实的。奥利金自己对两种意见都不满意。他说他避开了两种错误：我们这一派所坚持的关于肉身的错误，以及异端者所持的关于幻影的错误，因为双方都走向极端，有些人希望与从前一样，另一些人则完全否认身体的复活。他说：\Quote{有四元素，}他说，\Quote{为哲学家和医生所知：土、水、气、火，万物和人体都是由这些组合而成。我们在肉中看到土，在气息中看到气，在身体的湿气中看到水，在体温中看到火。因此，当灵魂奉天主的命令离开这易朽软弱的身体时，一切逐渐回归其母体物质：肉重新被吸收进土里，气息与空气混合，湿气回到深处，热量逸散到以太。就像你将一品脱的奶和酒倒入海中，又想再将混合之物分开一样，虽然你倒进去的奶和酒并未消失，但已不可能将倒出的东西分别保持；同样，血肉的本质就其原材料而言确实不消失，但它们不能再成为原来的结构，也不能完全与从前一样。}当他说这样的话时，他就否定了肉的坚实、血的流动、筋的致密、血管的缠绕和骨的坚硬。\par

26. \Quote{为另一个理由，}他说，\Quote{我们宣认我们身体的复活，就是那些已葬入坟墓并化为尘土的身体；保禄的身体将是保禄的，伯多禄的身体将是伯多禄的，每人将拥有自己的；因为灵魂在一个身体内犯罪而在另一个身体内受折磨，是不合理的，而且公义的审判者也不配使一个身体为基督流血，而另一个身体得冠冕。}谁听见这话，会认为他否认肉身的复活？\Quote{而且，}他说，\Quote{每一粒种子都因造物主天主的恩赐而内在地拥有其存在的法则，这法则以胚胎形式包含着未来的成长。粗壮的树木，连同其树干、树枝、果实、叶子，在种子中并未显现，但按潜在性或按希腊文所说的\InnerQuote{\OriginalTerm{种子理性}{spermatikos logos}}，它确实存在于种子中。在谷粒之内有一种髓质或脉管，当它在地里溶解时，吸引周围的物质，并以茎、叶和穗的形式重新升起；因此，当它死时是一样，当它从死者中复活时是另一样；因为在麦粒中，根、茎、叶、穗、干尚未分离。同样，在人的身体中，按照其存在的法则，某些原始要素存留以保证其复活，而一种髓质，即死者的苗床，在大地的怀抱中受滋养。但当审判之日来到，在总领天使的呼声和最后的号角声中，大地震动时，种子立即充满生命，并在顷刻间使死者迸发生命；然而它们所重新构成的肉身将不再是同样的肉身，也不会是原来的形状。为使你们确信我们说的是真理，让我引用宗徒的话：\BibleRef{格前 15:35}\InnerQuote{但有人会说：死人怎样复活呢？他们将带着什么样的身体来呢？糊涂人哪！你所播下的种子，若不先死，决不会生出生命来；并且你所播下的，并不是那将来要生的形体，而是一颗赤裸的种子，也许是麦子，或是葡萄树和树的种子。}我们已以麦粒和某种程度上以栽种树木为推理的对象，现在让我们以葡萄籽为例。它只是一粒小颗粒，小到几乎不能夹在两根手指之间。根在哪里？根、树干和枝条那缠绕的脉络在哪里？树叶的荫凉和那饱含未来美酒的可爱串串在哪里？你手指中的是干枯且几乎不可辨认的；然而，在那干枯的颗粒中，因天主的能力和繁殖的隐秘法则，那冒着泡沫的新酒必定有其起源。在树木上你允许这一切，在人身上你难道不允许这样的事可能发生吗？那朽坏的植物竟如此被装饰得美丽；我们为何认为那存留的人将恢复其先前的卑贱？你要求必须有肉身、骨头、血、四肢，以致你需要理发师来剪发，流鼻涕，剪指甲，下部生出污秽或服务于情欲？如果你引进这些愚蠢粗俗的念头，你就忘记了关于肉身的教导，即我们无法在其中中悦天主，且它是仇敌；你也忘记了关于死人复活的教导：\BibleRef{格前 15:42}\InnerQuote{播下的是可朽坏的，复活起来的是不可朽坏的；播下的是羞辱的，复活起来的是光荣的；播下的是软弱的，复活起来的是强壮的；播下的是属生灵的身体，复活起来的是属神的身体。}现在我们用眼观看，用耳听闻，用手行动，用脚行走。但在那属神的身体中，我们将全是视觉、全是听觉、全是行动、全是运动。主将改变\BibleRef{斐 3:21}我们卑贱的身体，使之相似祂光荣的身体。祂说\InnerQuote{改变}时，肯定了与我们现今所拥有的肢体的同一性。但一个\Quote{不同的}、属神的、属天的身体应许给我们，它既不可触摸，也不可目见，也不可衡量；它的变化将适应其未来居所的不同。否则，如果还是同样的肉身，如果我们的身体还是同样的，那么将再次有男有女，再次有婚姻；男人将有浓密的眉毛和飘动的胡须；女人将有光滑的脸颊和狭窄的胸膛，她们的身体必须适应怀孕和分娩。甚至幼小的婴儿也将复活；老人也将复活；前者需要哺乳，后者需要拐杖扶持。并且，愚昧的人哪，不要被我们主的复活所欺骗，因为祂显示了祂的肋旁和双手，站在岸上，与克罗帕一起行走，并说祂有肉有骨。那个身体，因为它不是出于男人的种子和肉欲，有其独特的特权。祂复活后吃了喝了，穿着衣服出现，并允许自己被触摸，为使祂怀疑的宗徒相信祂的复活。但祂仍然没有不显示一种属气和属神身体的本质。因为祂在门户关闭时进入，并在擘饼时消失不见。那么，难道我们在复活后也要吃要喝，并履行自然的职分吗？果真如此，那应许又成什么了？\BibleRef{格前 15:53}\InnerQuote{这可朽坏的必须穿上不可朽坏的。}}\par

27. 至此我们完全明白了，为何你在阐述信仰时，为了欺骗无知者的耳朵，九次提到身体，却一次也不提血肉，而众人还以为你承认血肉之身，以为血肉与身体是同一回事。如果它与身体是同一的，那就不表示任何区别。我这样说，是因为我知道你的回答：\Quote{我以为身体与血肉是相同的；我是极其单纯地说这话的。} 你为何不干脆称血肉为身体，随意地时而说血肉，时而说身体，从而表明身体由血肉构成，血肉即是身体呢？但请相信我，你的沉默并非单纯的沉默。因为血肉有一番定义，身体则有另一番定义；所有的血肉都是身体，但并非每个身体都是血肉。血肉特指由血液、血管、骨骼和筋腱组成的东西。尽管身体也被称为血肉，但有时它被描述为以太的或空气的，因为它不可触摸、不可看见；然而它又常常是可见且可触的。一堵墙是一个身体，但不是血肉；一块石头是一个身体，却不说它是血肉。因此，宗徒称有些身体是天上的，有些是地上的。天上的身体是指太阳、月亮、星辰的身体；地上的身体是指火、空气、水及其他，这些无生命的身体由物质元素构成。你看，我们识破了你的诡辩，并将你在内室中和在成全者中所讲的奥迹公之于众，这些奥迹本不该传入外人的耳中。你微笑着，举手弹指反驳说：\Quote{国王的女儿的一切荣耀都在内里。} 又引\BibleRef{歌 1:4}：\BibleQuote{国王引我进了他的内室。} 显然，你为何只谈身体的复活而不谈血肉的复活；当然是为了让我们这些无知者以为，当提到身体时，指的是血肉；然而成全者却明白，当提到身体时，血肉就被否定了。最后，宗徒在致哥罗森人书信中，为了表明基督的身体是由血肉做成的，而非属神的、属气的、稀薄的，便意味深长地说：\BibleRef{哥 1:21} \BibleQuote{连你们从前与基督疏远，在恶行上与祂的精神为敌，如今祂藉着死亡，在祂血肉的身体中使你们和好了。} 在同书信中又说：\BibleRef{哥 2:11} \BibleQuote{你们也在祂内，受了非人手所行的割损，是在脱去血肉的身体上，受了基督的割损。} 如果身体仅仅指血肉，且这个词没有歧义，也没有多种意义，那么同时使用\Emph{身体}和\Emph{血肉}这两个词就是完全多余的，好像身体不包含血肉似的。\par

28. 在我们信德与希望的信经中——这信经由宗徒传授，不是写在纸和墨上，而是写在肉质的心里——在宣认了天主圣三与教会的合一之后，整个基督徒教义的信经以血肉的复活结束。你却如此固执地专注于身体这个主题，在你的讲论中反复提起，先是身体，然后是身体，又是身体，九次提到身体，却一次也不提血肉；而宗徒们则总是谈论血肉，却只字不提身体。我要让你知道，我们看穿了你狡猾地添加的东西，以及你谨慎小心想要隐瞒的内容。因为你用与奥利金否认复活相同的章节来证明复活的真实性；你用模棱两可的论据来支持有疑问的立场，从而掀起一场风暴，顷刻间推翻了信仰的既定结构。你引用以下经文：\Quote{它被种下的是一个属生灵的身体，必要复活一个属神体的身体。} \Quote{因为人既不娶，也不嫁，却要像天上的天使一样。} 如果你是否认复活，你还会引别的例子吗？你说你打算真实而非想象地承认血肉的复活。在你用那些对无知者的耳朵加以粉饰的言辞，说我们带着死亡和埋葬时原有的身体复活之后，你为什么不再接着说：\Quote{主在复活后，显明了祂手上的钉痕，指出了祂肋旁的枪伤，当宗徒们怀疑，以为看见的是幽灵时，祂回答他们说：\BibleRef{路 24:39} \InnerQuote{你们摸摸我，看看！因为幽灵没有血肉和骨头，如同你们看到我有的。} 特别对多默说：\BibleRef{若 20:27} \InnerQuote{把你的指头伸到我的手里来，把你的手伸到我的肋旁来，不要作无信的人，但要作有信德的人。} 同样，在复活后，我们也会拥有我们现今使用的同样肢体，同样的血肉、血液和骨骼，因为圣经所谴责的不是这些的本质，而是它们的行为。此外，在创世纪中记载：\BibleRef{创 6:3} \InnerQuote{我的神不会常住在那些人里面，因为他们是血肉。} 宗徒保禄论及犹太人的败坏教义和行为时说：\BibleRef{迦 1:16} \InnerQuote{我没有同血肉之人商量。} 他又对显然在血肉中的圣人们说：\BibleRef{罗 8:9} \InnerQuote{你们却不在血肉内，而在圣神内，因为天主的圣神住在你们内。} 他否认那些明明在血肉中的人是在血肉内，并不是谴责血肉的实体，而是谴责它的罪过。}\par

29. 关于复活的正确宣认宣告：肉身将是光荣的，但并不毁灭其真实性。当宗徒说：\BibleQuote{这可朽坏、可死的}\BibleRef{格前 15:53}，他的话是指\Emph{这同一身体}，即当时所见的肉身。但当他接着说它穿上不朽和不死时，他并不是说那穿上之物——即那件衣裳——除去了它所装饰在光荣中的身体，而是使那先前缺乏光荣的身体成为光荣的。如此，抛弃了可死和软弱这更无价值的衣裳，我们便可穿戴上不朽的金饰，并且，可以说，也穿戴上力量与德能的福乐；因为我们不愿被剥去肉身，而是要在其上穿上光荣的外衣，渴望穿上那自天而来的居所，好使可死的被生命吞灭。毫无疑问，没有人能穿上外衣，除非他先前已穿着衣服。同样，我们的主在山上是这样显圣容的：祂并未失去双手、双脚或其他肢体，突然开始像太阳或球体那样滚动成圆形；而是同一肢体闪耀着太阳般的光辉，使宗徒们的眼睛失明。因此，祂的衣服也改变了，但只是变得洁白发光，并非成为气态的；因为我想你并不打算坚持说祂的衣服也是属神的。\BibleRef{玛 17:2} 圣史接着说祂的面貌发光如太阳；但谈及祂的面貌时，我认为祂的其他肢体也同样被看见了。\Person{哈诺客}带着身体被提升；\Person{厄里亚}带着身体被提往天上。他们并未死去，而是乐园的居民；即使在乐园中，他们仍保有那被提去和迁移时所拥有的肢体。我们藉着禁食所追求的，他们却因与天主的共融而拥有。他们以天上的食粮为生，并以天主的每一句话为满足，以那同为他们的主的天主为食物。请听救主说：\Quote{并且我的肉身要安息于希望中。}在别处又说：\BibleQuote{祂的肉身并未见到腐朽。}\BibleRef{宗 2:31}以及：\BibleQuote{凡有血肉的，都要看见天主的救恩。}\BibleRef{依 40:5}难道你要永远把身体说成是双重的东西吗？还是请引用\Person{厄则克耳}的异象：他将骨头与骨头连接起来，从坟墓中带出，然后使它们站立起来，用肉和筋联结，用皮包裹。\par

30. 请听\Person{约伯}——苦难的征服者——所说如同雷鸣的话；他正用瓦片刮去腐坏肉身的污秽，用复活的希望和现实来安慰自己的痛苦：\Quote{哦，但愿}他说，\Quote{我的言语都记录在册！哦，但愿它们被铁笔写在书上，刻在铅版上，凿在岩石上，永远存留！因为我知道我的救赎主活着，末日我必从地上起来，再次以我的皮包裹我，并在我的肉身内看见天主，我要亲眼看见祂，我的眼睛要注视，并非别人。这希望存于我怀中。}还有什么比这预言更清楚的呢？自基督以来，没有人在祂之前像祂那样公开谈论复活。祂希望自己的话永远流传；为使它们不被岁月磨灭，祂愿写在铅版上，刻在岩石上。祂盼望复活；不，祂更是知道并看见祂的救赎主基督活着，并将在末日从地上复活。主那时尚未死去，而教会中的这位斗士却已看见祂的救赎主从坟墓中复活。当祂说：\Quote{我再次以我的皮包裹我，并在我的肉身内看见天主}，我想祂这样说并非出于对肉身的喜爱，因为那肉身正在祂眼前腐烂发臭；而是出于复活的信心，并因着未来的安慰，祂轻看眼前的痛苦。祂又说：\Quote{我必以我的皮包裹我。}这里我们哪里找到属天的身体？哪里找到如同气息和风一样的属气的身体？在有皮和肉，有骨头、筋、血和血管的地方，必定是血肉的构造和性别的区分。祂说：\Quote{在我的肉身内，我要看见天主。}当凡有血肉的都看见天主的救恩，并看见耶稣就是天主时，那么我也要看见救赎主、救主和我的天主。但我将在那现在折磨我、因痛苦而消融的肉身中看见祂。因此，在我的肉身内我要看见天主，因为祂藉自己的复活医治了我一切的软弱。你以为\Person{约伯}那时是在反对\Person{奥力振}，并为了肉身的真实性——在这肉身中他忍受了折磨——而与异端者进行与我们相似的争论吗？因为祂不愿认为自己的一切痛苦都是徒劳的；而实际上祂所承受的肉身既作为肉身受了折磨，那么复活的将会是另一种属神的肉身。因此祂强调并坚持真理，并且通过直截了当地说出以下的话，阻止了一切可能隐藏在狡诈的宣认中的东西：\Quote{我要亲眼看见祂，我的眼睛要注视，并非别人。}如果祂不在自己的性别中复活，如果祂不拥有那时躺在粪堆上的同一肢体，如果祂不睁开那当时正在看虫子的同一眼睛来看天主，那么\Person{约伯}在哪里？你取消了构成\Person{约伯}的一切，却给我空洞的说法\Emph{约伯将要复活}；这就好像你命令一艘船在沉船后复原，却拒绝提供构成船的任何具体东西。\par

31. 我要坦白地说，虽然你们扭曲嘴巴、扯头发、跺脚，像犹太人一样拿起石头，我仍要公开宣认教会的信仰。没有血肉、没有骨骼、没有血液、没有肢体的复活，是不可理解的。哪里有血肉、骨骼、血液和肢体，就必然有性别的差异。哪里有性别的差异，那里若望就是若望，玛利亚就是玛利亚。你们不必担心那些在死前曾经生活在自己的性别中而没有履行性别功能之人的婚姻。当经上说：\Quote{在那一天，他们不娶也不嫁}，这话是指那些可以结婚却不会结婚的人。因为没有人论及天使说：\Quote{他们不娶也不嫁}。我从未听说过天堂的精神德性中举行婚礼：但哪里有性别，那里就有男人和女人。因此，虽然你们不情愿，却被真理迫使承认：\Quote{人要么因在身体内度纯洁正直的生活而得到冠冕，要么因在身体内作了享乐与邪恶的奴隶而受到惩罚。}以\Emph{肉体}替换\Emph{身体}，你并没有否认男女的存在。如果我们没有使不洁成为可能的性别，谁能因守贞的生活而获得荣耀呢？谁曾因一块石头保持童贞而给它加冕呢？应许给我们的是与天使相似，也就是说，我们将在我们的肉体和性别中，得到他们那没有肉体、没有性别的天使存在的福乐。我如此单纯地相信，并知道如此宣认：性别可以在没有感官功能的情况下存在；人就是这样复活，并且是这样与天使同等。肢体的复活也并非立即显得多余，因为它们将没有功用，因为仍在今生时，我们已努力不去履行肢体的工作。此外，与天使相似并不意味人变成天使，而是他们在不朽与光荣中的成长。\par

32. 至于你用来反对教会的那些从孩童、婴儿、老人、食物和排泄物中得出的论证，其实并非你自己的；它们来自外教人的源头。因为外教人用同样的论点嘲笑我们。你说你是基督徒；放下外教人的武器。他们应该从你那里学习宣认死人的复活，而不是你从他们那里学习否认复活。或者，如果你属于敌人的阵营，就公开表明你是敌对者，这样你就可以分担我们加给外教人的创伤。我允许你嘲笑说需要奶妈来阻止婴儿哭泣；以及那些衰老的老人，你担心他们会在冬天的寒冷中干瘪。我也承认，理发师学艺毫无用处，因为我们岂不知道以色列子民四十年间指甲和头发都没有生长；更有甚者，他们的衣服没有穿破，鞋子也没有穿旧？\Person{哈诺客}和\Person{厄里亚}，我们刚才还在谈论他们，至今仍停留在被带走时的状态。他们有牙齿、肚子、生殖器官，却不需要食物或妻子。你为什么诽谤天主的能力？祂能从你所提到的\Emph{精髓}和\Emph{种子之地}中，不仅从肉体产生肉体，还能从一个身体造出另一个身体；并且将水——即无价值的肉体，变成灵性身体的宝贵美酒。祂从无中创造万物的同一能力，也能恢复已经存在的东西，因为修复已有的东西比创造从未有过的东西要小得多。看到一个人从泥土中被创造出来，没有经历成长的各个阶段，你怎能对从婴儿和老年状态复活到成熟状态感到惊奇呢？一根肋骨变成了女人；通过创造人的第三种方式，我们出生时那些令我们羞耻的卑劣元素变成了肉，由肢体连在一起，流成血管，硬化成骨骼。还有第四种人类的繁衍方式，我可以告诉你：\BibleQuote{圣神要降临于你，至高者的能力要庇荫你，因此，那圣者将称为天之子。}\BibleRef{路 1:35}\Person{亚当}以一种方式被造，\Person{厄娃}以另一种，\Person{亚伯尔}以另一种，那人\Person{耶稣基督}又以另一种。然而，所有这些开端虽然不同，人的本性却始终如一。\par

33. 若我欲证明肉身及所有肢体的复活，并解释数段经文的含义，便需多卷著作；但当前之事并不要求如此。因为我本意不是要逐条答复\Person{奥力振}，而是要揭露你那不诚恳的\Quote{\WorkTitle{辩护辞}}的奥秘。然而，我在反驳你的立场上已停留过久，并担心我急于揭露欺骗，反而会给读者留下绊脚石。因此，我将汇集证据，并在行经时略略一瞥这些证明，好让我们能用圣经的全部重量来攻击你那有毒的论点。那没有婚宴礼服、又没有遵守\BibleRef{训 9:8} \BibleQuote{你的衣裳要时常洁白，}这一命令的，手脚被捆绑，不得在宴席上斜靠，或坐在宝座上，或站在天主的右边；\BibleRef{玛 22:13} 他被丢到地狱，那里有哀号切齿。\BibleRef{路 12:7} \BibleQuote{你们的头发都被数过。}如果头发都被数过，我想牙齿就更容易被数了。但如果它们终有一日要毁灭，数它们就毫无意义。\BibleQuote{时候将到，凡在坟墓里的都要听见天主子的声音，并要出来。}他们将用耳朵听见，用脚出来。\Person{拉匝禄}已经这样做了。而且，他们要从坟墓里出来；也就是说，那些被安置在坟墓里的死者要来，并从他们的墓穴中复活。因为天主所赐的甘露是医治他们骨头的。那时，天主藉先知所说的话将应验：\BibleRef{依 26:20} \BibleQuote{我的百姓，去进入你的内室，关闭你的门，隐藏片刻，直到忿怒过去。}内室代表坟墓，当然，被放进坟墓的都要从坟墓里出来。他们要像脱缰的幼骡从坟墓中出来。他们的心要欢喜，骨头要像太阳般升起；所有肉身都要来到主面前，祂要命令海中的鱼；它们要交出所吞的骨头；祂要将关节对关节，骨对骨；\BibleRef{达 12:2} \BibleQuote{那些睡在尘土中的必醒来，有些进入永生，有些蒙受羞辱和永远可耻。}那时，义人要看见恶人的刑罚和折磨，因为\BibleRef{依 66:24} \BibleQuote{他们的虫子不死，他们的火不灭，他们必被所有肉身看见。}因此，我们凡有这希望的，既然我们曾将肢体献给不洁和不法以行不法，如今也要将肢体献给正义以行圣化，好使我们能\BibleRef{罗 6:4} \BibleQuote{从死者中复活，并在新生活中行走。}同样，主耶稣的生命彰显在我们必死的身体上，\BibleRef{罗 8:11} \BibleQuote{那使耶稣基督从死者中复活的，也必要藉着住在我们内的圣神，使我们必死的身体活起来。}因为这是应当的：我们既常将基督的致死带在身上，耶稣的生命也应当彰显在我们必死的身体上，即我们的肉身上；这肉身按本性是必死的，但按恩宠却是永恒的。\Person{斯德望}也看见耶稣站在圣父的右边，\BibleRef{出 4:6} \BibleQuote{梅瑟的手变得雪白，后来又恢复原色。}这仍然是手，尽管两种状态不同。\Person{耶肋米亚}中的陶工，他所做的器皿因石头的粗糙而破碎，他仍用同一团泥、同一黏土恢复了破碎的器皿；并且，如果我们看\Emph{复活}这个词本身，它并不意味着一个东西被毁灭，另一个被升起；而且加上\Emph{死亡}这个词，正是指向我们的肉身，因为人身上死亡的部分，也正是被赋予生命的。\BibleRef{路 10:34} 在往\Place{耶里哥}的路上受伤的人，被带到客栈，肢体完整，他罪过的伤痕因不死而痊愈。\par

34. 甚至坟墓也开了\BibleRef{玛 27:52}，在我们主受难时，太阳躲避，大地震动，许多圣人的身体起来，在圣城中显现。\BibleQuote{这是谁，}依撒意亚说，\BibleQuote{那从\OriginalTerm{厄东}{Edom}上来，从\OriginalTerm{波斯拉}{Bozrah}而来，衣裳闪耀，如此荣美？}\OriginalTerm{厄东}{Edom}按解释，是\Emph{属土的}或\Emph{血腥的}；\OriginalTerm{波索}{Bosor}是\Emph{肉体}或\Emph{在患难中}。他用几句话揭示了复活的全奥秘，即肉体的真实性和荣耀的增长。含义是：那从土里上来、从血里上来的是谁？按照\BibleRef{创 49:11}雅各伯的预言，他把自己的驴驹拴在葡萄树上，独自踹了酒榨，他的衣服因从\OriginalTerm{波索}{Bosor}来的新酒而变红，即从肉体，或从世界的患难而来：因为\BibleRef{若 16:33}他自己已战胜了世界。因此，他的衣服是红色而闪耀的，因为他比世人更俊美，因他胜利的荣耀，它们变成了白袍；然后，关于基督的肉体，确实应验了\BibleRef{歌 8:5}的话：\BibleQuote{那全身洁白，靠着她爱人的是谁？}以及同一书卷中所写的\BibleRef{歌 5:10}：\BibleQuote{我的爱人洁白而红润。}这些人是他的真门徒，他们没有\BibleRef{默 14:4}与女人玷污衣服，因为他们保持了童贞，为天国自阉。所以他们要穿白衣。然后，我们主的话将完全应验：\BibleRef{若 6:39}\BibleQuote{凡父赐给我的，我必不失落其中任何一部分，反而要在末日使他复活；}他的整个人性，即他出生时完全取来的。然后，那失落、在阴间游荡的羊，将被救主扛在肩上完整地带回；那因罪而病的羊，将受到审判者仁慈的扶持。然后，他们将看见那刺透他的人，那曾喊叫\BibleRef{若 19:6}\BibleQuote{钉死他，钉死他}的人。他们将一次又一次捶胸，他们和他们的女人，就是那些我们主背着十字架时所说\BibleRef{路 23:28}\BibleQuote{耶路撒冷的女子，不要为我哭，当为你们自己和你们的儿女哭}的那些女子。然后，天使对惊愕的宗徒所说的预言将应验：\BibleRef{宗 1:11}\BibleQuote{加里肋亚人，你们为什么站着望天？这位离开你们被接到天上去的耶稣，你们见他怎样升了天，也要怎样降来。}但我们对一个这样说的人该作何感想：他说我们的主复活后与宗徒同食四十天，以免他们以为他是幻影，然后又断言正是这幻影做了这事，吃了，并被人亲眼看见。所看见的，要么是真实的，要么是虚假的。如果是真实的，那么他确实吃了，确实有肢体。如果是虚假的，他怎能愿意用虚假的印象来证明他复活的真实性？因为没有人用虚假来证明真实。你会说，那么我们复活后也要吃吗？我不知道。圣经没有告诉我们；然而，若问这个问题，我认为我们不会吃。因为我读到天主的国不在于吃喝，它却应许了\BibleRef{格前 2:9}\BibleQuote{眼睛未曾看见，耳朵未曾听见，人心也未曾想到的事}。梅瑟禁食了四十昼夜。人性不允许这样，但在人不可能的，在天主却不然。正如预言未来时，一个人宣布十年后还是一百年后发生的事并不重要，因为对未来的知识是同一的；同样，一个人能禁食四十天仍活着——不是说他靠自己禁食，而是靠天主的权能活着——也将能永远不饮不食而活着。为什么我们的主吃了一个蜜巢？为了证明复活：不是为了让你的味觉享受蜂蜜的滋味。他要求一条烤在炭火上的鱼，以便\BibleRef{若 21:9}坚定那些怀疑的宗徒，他们不敢靠近他，因为他们以为自己看见的不是身体，而是神魂。\EditorialNote{马尔谷福音第五章}会堂长的女儿复活后也吃了食物。\EditorialNote{若望福音第十二章}死了四天的拉匝禄复活，并在宴席上出现在我们面前；不是因为他在阴间习惯吃东西，而是因为一个如此困难的事例考验了信徒的批评。正如他向他们显示了真实的手和真实的肋旁，他也真实地与门徒一同吃了；真实地与\Person{克罗帕}同行；用真实的舌头与人交谈；真实地斜靠在晚餐上；用真实的手拿起饼，祝福了，掰开，递给他们。至于他突然从他们眼前消失，那是天主的权能，而非虚幻的幻影。此外，甚至在他复活之前，当他们带他出\Place{纳匝肋}，想把他从山崖推下去时，他从他们中间穿过去，即从他们手中逃脱了。我们能跟随\Person{马尔西翁}，说因为他被抓住时以一种违反本性的方式逃脱了，所以他的诞生只是虚假的吗？难道主没有魔法师也享有的特权吗？据说\Person{提亚纳的阿波罗尼乌斯}在\Person{多米提安}面前站在法庭上时，突然消失了。不要将主的权能与魔法师的把戏等量齐观，以致他显得是而非是，被人以为他无牙而食，无足而行，无手而掰饼，无舌而说话，并显示了一副没有肋骨的肋旁。\par

35. 你可能会说，如果祂拥有和先前一样的身体，为什么他们在路上没有认出祂？让我回想圣经所说：\BibleRef{路 24:16} \BibleQuote{他们的眼睛被蒙蔽了，以致认不出祂。} 又说：\BibleQuote{他们的眼睛开了，就认出了祂。} 当祂不被认出时是一人，被认出时是另一人吗？祂当然始终如一。因此，他们认识与否，取决于他们的视觉；不取决于被看见的祂；然而在某种意义上也取决于祂，因为祂蒙蔽了他们的眼睛，使他们认不出祂。最后，为了让你明白，迷惑他们的错误不应归咎于主的身体，而应归咎于他们的眼睛被遮蔽，我们被告知：\EditorialNote{若望福音二十章} \BibleQuote{他们的眼睛开了，就认出了祂。} 所以，\Person{玛利亚玛达肋纳}在尚未认出耶稣，并在死人中寻找活人时，以为祂是园丁。后来她认出了祂，便称祂为主。复活后，耶稣站在岸上，门徒们在船中。当其他人没有认出祂时，耶稣所爱的门徒 \BibleRef{若 21:7} 对\Person{伯多禄}说：\BibleQuote{是主。} 因为贞洁最先认出贞洁的身体。\Quote{祂是同一的，但并非所有人都同样地看见祂是同一的。} 紧接着又记载：\BibleQuote{没有人敢问祂：你是谁？因为他们知道祂是主。} 没有人敢问，因为他们知道祂是天主。他们与祂一同用晚餐，因为他们看见祂是有人性、有肉身的人；并非祂作为天主是一人，作为人是另一人，而是，祂是同一天主子，作为人而被认识，作为天主而被朝拜。我想我现在必须卖弄一下哲学，说我们的感官不可靠，尤其是视觉。必须从死者中唤醒一位\Person{卡内亚德}来告诉我们真相——船桨在水里看似折断，远处的柱廊看起来更宏伟，塔楼的尖角在远处显得圆润，鸽子的背部随移动而变色。当\Person{洛达} \EditorialNote{宗徒大事录十二章} 宣布\Person{伯多禄}到来，并告诉宗徒们时，由于危险巨大，他们不相信他已逃脱，反而怀疑是个幻影。此外，祂穿过关闭的门，显示了与从视线中消失同样的能力。正如传说中，\Person{林叩斯}能透过墙壁看见。难道主不能穿过关闭的门进入，除非祂是幻影？鹰和雕能跨海看见死尸。难道救主不能不开门就看见祂的宗徒？告诉我，最敏锐的辩论者，哪一样更伟大：将地球的巨大重量悬于虚无，并在变化的波面上保持平衡；还是天主穿过关闭的门，让受造物屈服于造物主？你承认了更大的，却反对较小的。\Person{伯多禄} \BibleRef{玛 14:28} 以他沉重而坚实的身体在水面上行走。柔软的水并没有屈服：他的信德稍有动摇，身体立刻明了自身的本性；让我们知道，在水面上行走的不是他的身体，而是他的信德。\par

36. 我请求你们，这些用如此精巧的论点反对复活的人，让我们简单谈一谈。你相信我们的主真的以祂死去并被埋葬的同一身体复活了，还是不相信？如果你相信，为什么你提出那些导致否认复活的命题？如果你不相信，你就是这样试图欺骗无知者的思想，虽然有口无心地炫耀\Quote{复活}这个词，那么听我说。不久前，有一位\Person{马尔西昂}的门徒说：\Quote{那带着这血肉和这骨头复活的人有祸了！} 我们的心立刻以喜乐回应，\BibleRef{罗 6:4} \BibleQuote{我们藉着洗礼已与基督同葬，也将与祂一同复活。} \Quote{你指的是灵魂的复活，还是肉身的复活？} 我回答说：\Quote{不仅是灵魂的，也是肉身的，它连同灵魂一起在洗礼池中重生。凡在基督内重生的，怎能朽坏呢？} \Quote{因为经上记载，} 他说，\BibleRef{格前15:50} \Quote{\InnerQuote{血肉之躯不能承继天主的国。}} \Quote{我恳请你注意所说的话——\InnerQuote{血肉之躯不能承继天主的国。}} \Quote{这是说他们不会复活。} \Quote{根本不是，只是说\InnerQuote{他们不能承继那个国。}} \Quote{怎么回事？} \Quote{\InnerQuote{因为，}接着，\BibleRef{格前15:50} \InnerQuote{朽坏的也不能承继不朽坏的。}因此，只要他们仍是纯粹的血肉之躯，就不能承继天主的国。但当可朽坏的穿上了不朽坏的，必死的穿上了不死的，肉身的泥土被制成器皿时，那先前被沉重重量压在地上的肉身，一旦接受了圣神的翅膀——这翅膀意味着它的改变，而非毁灭——将以新的荣耀飞向天堂；那时将应验经上所写的：\InnerQuote{死亡被胜利吞噬了。死亡啊，你的夸耀在哪里？死亡啊，你的刺在哪里？}}\par

37. 我们颠倒次序，回答了关于灵魂状况和肉身复活的问题；并省略了信函的开头部分，只限于反驳这篇非常出色的论著。因为我们宁愿谈论天主的事，而不愿谈论我们自己的冤屈。\BibleRef{撒上 2:25}\BibleQuote{若一人得罪了另一人，他们可以为他向主祈祷。但若他得罪了天主，谁能为他祈祷呢？}如今，与此相反，我们将首要之事放在以不灭的仇恨追捕那些伤害我们的人——而对于那些亵渎天主的人，我们却宽厚地伸出手。\Person{若望}向\Person{德奥斐洛}主教写了一篇辩护词，其开头这样写道：\Quote{你，确实，作为天主的人，以宗徒的恩宠为装饰，肩负着对所有教会的关怀，尤其是\Place{耶路撒冷}的教会，尽管你自己因你治下天主的教会而分心于无数的忧虑。}这是赤裸裸的谄媚，企图将权力集中在个人手中。你，一方面要求教会的规则，并运用尼西亚大公会议的法令，并且声称对属于另一个教区、实际上与自己的主教同住的圣职人员拥有权威，回答我的问题：\Place{巴勒斯坦}与\Place{亚历山大里亚}主教有什么关系？除非我受骗，那些法令规定\Place{凯撒勒雅}是\Place{巴勒斯坦}的大都会，\Place{安提约基}是整个\Place{东方}的大都会。因此，你应当要么向\Place{凯撒勒雅}主教申诉，你知道我们与他共融，而拒绝与你共融；要么，如果需要在远处寻求审判，书信应当寄往\Place{安提约基}。但我知道你为什么不愿意寄往\Place{凯撒勒雅}或\Place{安提约基}。你知道要逃避什么、避免什么。你宁愿用你的抱怨去攻击那些先入为主的耳朵，而不是向你的大主教致以应有的敬意。我这样说，并不是因为我在这次使命本身有什么可指责的，除了某些引起怀疑的偏袒，而是因为你应当在实际质问者们面前为自己辩白。你以这些话开头：\Quote{你派遣了一位最忠诚的天主仆人，司铎\Person{依西多尔}，一位无论从举止和服饰的尊严，还是从神圣的理解力方面都颇具影响力的人，去医治那些灵魂重病的人；但愿他们能意识到自己的病！一位天主的人派遣一位天主的人。}司铎与主教之间没有区别；派遣者与受遣者享有同样的尊严；这够跛脚的了；正如俗语所说，船在港口搁浅了。那位你用赞美捧上天的\Person{依西多尔}，在\Place{亚历山大里亚}与你同在\Place{耶路撒冷}一样，背负着同样的异端指控；因此，他到你这里来似乎不是作为特使，而是作为同谋。此外，在他亲笔写的信函中——这些信函在派遣使团前三个月因地址错误而寄给了我们——被交给了司铎\Person{文森修斯}，至今仍由他保管。在这些信中，写信人鼓励他军队的领袖将脚踏稳在信德的磐石上，不要因我们的哀歌而惊恐。他在我们对其使命有任何怀疑之前就承诺，他将来\Place{耶路撒冷}，他一到，他对手的阵营将立即被粉碎。除此之外，他还用了这些话：\Quote{正如烟雾消散在空中，蜡在火旁熔化，那些永远反对教会信德、现今藉着简单人试图扰乱这信德的人，也将这样四散。}\par

38. 我问你，我的读者，一个人在他到来之前写这些东西，在你看来，他是什么？一个对手，还是一个特使？这个人，我们确实可以称他为最虔诚的，或最热心的，并且，给出这个词的准确对等词，一个献身于崇拜天主的人。这就是那位拥有神圣理解力的人，如此有影响力，在举止和服饰上如此尊严，以至于，像一位属灵的\Person{希波克拉底}，他能藉着临在缓解我们灵魂的疾病，只要我们愿意接受他的治疗。如果这就是他的药，让他治愈自己吧，因为他习惯于治愈他人。对我们来说，他那神圣的理解力是为\Person{基督}的缘故而愚拙。我们宁愿停留在我们简单的疾病中，也不愿通过使用你的眼药，学会一种不虔诚的滥用视觉。接下来是这些话：\Quote{您圣善的卓越意图迫使我们日夜向主祈祷；并且，好像这些意图已经完美实现一样，我们在圣地献上我们的祈祷，求祂赐给您完美的赏报，并赐给您生命的冠冕。}你感谢得好；因为，如果\Person{依西多尔}没有来，你现在在整个\Place{巴勒斯坦}找不到如此忠实的伙伴。如果他没有带来他预先承诺的帮助，你会发现自己被一群无法理解你智慧的乡巴佬包围。现在正在讨论的这篇辩护词就是在\Person{依西多尔}的在场下、并在很大程度上在\Person{依西多尔}的协助下口授的，以至于同一个人既撰写了信函又将其送达目的地。\par

39. 您的信继续叙述：\InnerQuote{他来到这里，与我们进行了三次单独的会晤，并用您神圣的智慧和自己的理解力所发出的治愈性语言处理此事，却发现他对任何人都没有帮助，也没有人能帮助他。}事实是，那位据说\InnerQuote{与我们进行了三次单独的会晤}、以期在他的到来中保持神秘数字，并与我们谈论了由\Person{德奥斐洛}主教发布的命令的人，不愿交给我们他带给我们的信件。当我们说：\Quote{如果你是使者，请出示你的凭证；如果你没有信件，如何向我们证明你是使者？}他回答说，他确实有给我们的信件，但他被耶路撒冷主教要求不要交给我们。你看，这位真正的使者与他自己的身份一致；你看他如何对双方都表现出公正，以促成和平，并排除偏袒任何一方的嫌疑。无论如何，他没有带着药膏而来，也没有掌握医生的工具，因此他的药物无效。您接着说：\InnerQuote{\Person{热罗尼莫}和与他一起的人，既在暗中，也在众人面前，再三并起誓，使他确信他们从未对我们的正统信仰有过怀疑，说：\Quote{我们现在对他在信仰方面的感觉，与过去我们与他共融时完全相同。}}看，教义上的一致能做什么。\Person{依西多禄}为了能做出这样的报告，被接纳为亲密伙伴，被称为天主的人，最虔诚的司铎，一个有影响力的人，神圣可敬的举止和神圣的理解，基督徒的\OriginalTerm{希波克拉底}{Hippocrates}。我这个可怜虫，躲在隐居中，突然被这位伟大的主教隔绝，失去了司铎的名号。那么，这个\Person{热罗尼莫}，带着他衣衫褴褛的牧群和卑微的追随者，还敢对\Person{依西多禄}和他的霹雳做出任何回答吗？当然不敢；无非是害怕这位使者绝不会让步，并且他的出现和巨大身材可能压倒他们。\InnerQuote{不是一次，也不是三次，而是再三起誓，他们知道那人正统，从未怀疑过他是异端者。}多么露骨和无耻的谎言！一个人为自己作的见证！这样的见证即使在\Person{卡托}口中也不被相信，因为凭两三个见证人的口，每句话才得确立。可曾有哪句话或哪个口信传给您，说我们在未确信您正统的情况下，会容忍与您共融？当通过最杰出的基督徒\Person{阿尔赫劳斯}伯爵，他试图在我们之间促成和平时，曾指定了一个我们会面的地点，难道首要条件之一不就是信仰应成为未来协议的基础吗？他答应来。逾越节临近；大批隐修士聚集；您应在约定的地点出现；您不知该做什么。突然您传话说某人生病，那天不能来。这是演员还是主教在说话？假设您说的是真的，为了迎合一个软弱女人的喜好，她担心您不在时可能会头痛、恶心或胃痛，您就忽视教会的利益吗？您蔑视这么多聚集的基督徒和隐修士吗？我们不愿给中断谈判提供借口；我们看穿了您拖延的诡计，试图以忍耐克服您对我们的不公。\Person{阿尔赫劳斯}再次写信，建议他他会多待两天，以防他愿意来。但他很忙；他亲爱的女人还在呕吐，他无暇顾及我们，直到她摆脱恶心。两个月后，期待已久的\Person{依西多禄}终于到达，他从我们这里听到的不是如您所声称的为您作证，而是我们为何要求他澄清的理由。因为当他提出：\InnerQuote{如果他是异端者，你们为何与他共融？}我们所有人都回答，我们是在毫无怀疑他是异端者的情况下与他共融的；但在他被最可敬的\Person{厄丕法尼}以言语和信件传唤，却拒绝回答之后，\Person{厄丕法尼}本人致函隐修士们，除非他就信仰做出澄清，否则不可轻率地与他共融。信件在我们手中；此事确凿无疑。因此，这是全体弟兄的回答：并非如您所坚持的那样，您不是异端者，因为过去您没有被说成是异端者。因为按照那种说法，一个人之前健康，就不能说他生病了。\par

40. 继续看信。\InnerQuote{但当\Person{保利尼亚诺}及与他一起的人的圣秩被提出时，他们开始感到自己理亏。为了爱德与和谐，对他们做出了所有让步，唯一坚持的是，尽管他们违背规则被祝圣，但仍应服从天主的教会的权威，不应分裂教会，建立自己的权威。但他们不同意这一点，开始提出关于信仰的问题；这样，他们就向所有人表明，如果\Person{热罗尼莫}司铎和他的朋友们没有被指控，他们就对我们没有指控，而他们之所以诉诸教义问题，是因为当他们被指控错误和不当行为时，他们完全无法就这类事情回答我们，也无法对他们的不当行为给出任何满意的解释：并非他们希望我们能因异端被定罪，而是他们试图损害我们的声誉。}\par

41. 无人应因这冗词而责怪译者：希腊文原文即是如此。同时，我欣喜的是，尽管我以为自己已被斩首，却发现我的司铎头又回到了肩上。他说我们完全无法被说服，且他回避交锋。如果分歧的原因不在于关于信仰的讨论，而是源于保利尼亚诺的祝圣，那么拒绝回答以给那些寻找借口者以借口，岂非极端的愚蠢？宣认信仰；但务要如此回答所向你提出的问题，使众人明白争论不是关于信仰，而是关于秩序。因为只要你在被问及信仰时保持沉默，你的对手就有权对你说：\Quote{这问题关乎信仰，而非秩序。}如果是秩序问题，你在被问及信仰时一言不发就是愚行。如果是信仰问题，你以秩序问题为借口就是愚行。此外，当你说你的目的是要他们服从教会、不分裂教会、不建立自己的权威时，你所指的人是谁，我不太明白。如果你指的是我和司铎文森提，那么你睡得太久了，如果直到十三年后的现在才醒来说出这些话。因为我离开安提约基雅、他离开君士坦丁堡（两座名城）的原因，并非要赞美你受欢迎的雄辩，而是要在乡间和隐居中为我们青年时期的罪过哭泣，并为我们赢得基督的怜悯。但如果你所论的是保利尼亚诺，他正如你所见，服从他的主教，住在塞浦路斯：他有时来看望我们，并非作为你的神职人员，而是作为别人的，即祝圣他的那个人的。但即使他想留在这里，过一种安静、独居的生活，分享我们的流亡，他对你有什么亏欠，除了我们该对所有主教表示的尊敬？假定他曾被你祝圣，他只会对你说同样的话，如同我这个可怜人对已故的保禄主教所说的：\Quote{难道我曾请求被你祝圣吗？}我说，\Quote{倘若在授予司铎品秩时不剥夺我们隐修的生活，你可以按你所认为最好的方式授予或保留祝圣。但倘若你授予司铎之名的意图是夺去我为此而舍弃世界的那个事物，我仍必须坚持我一贯所是的；你因祝圣我并未遭受任何损失。}\par

42. \Quote{使他们不分裂教会，}他说，\Quote{并建立自己的权威。}谁分裂教会？是我们，作为白冷的一个完整家庭在教会中共融？还是你，要么是正统的信徒却因骄傲拒绝谈论信仰而沉默，要么是异端而真正分裂教会？是我们分裂教会吗？几个月前，大约在五旬节，当日头变黑、全世界都畏惧审判者即将来临时，我们向你的司铎提交了四十位不同年龄和性别的候洗者接受圣洗。在隐修院中确实有五位司铎有权付洗，但他们不愿做任何可能激怒你的事，以免你以此为借口对信仰保持缄默。相反，难道不是你分裂教会吗？你命令你在白冷的司铎在复活节不要给我们的人付洗，以致我们不得不送他们去第奥斯波利斯，由宣信者兼主教狄约尼削付洗。我们说我们分裂教会吗？我们在修道室之外在教会中不担任任何职务。难道不是你反而分裂教会吗？你向你的神职人员颁布命令：若有人说保利尼亚诺是由埃皮法尼祝圣为司铎的，就禁止他进入教堂。从那时起直到今日，我们只能从外面注视救主的洞穴，而异端者却进入，我们则站在远处叹息。\par

43. 我们是分裂者吗？难道不是那个拒绝给活人以居所、给死人以坟墓、并要求流放自己弟兄的人才是分裂者吗？是谁特别疯狂地把那头不断威胁全世界咽喉的野兽放到我们脖子上？是谁让雨水直到此刻仍落在圣人的骸骨和他们无害的灰烬上？这些就是善牧用来邀请我们和好的亲爱之举，同时却指控我们建立自己的权威——我们这些在共融和爱德中与所有主教联合的人，至少在他们为正统之时。难道你自己就是教会，而凡得罪你的就被摒于基督之外？如果我们捍卫我们自己的权威——证明我们在你的教区有一位主教。我们未曾与你共融的原因在于信仰问题；回答我们的问题，它就会变成秩序问题。\par

44. \Quote{他们，}你接着说，\Quote{也利用他们声称\Person{厄丕法尼}写给他们的其他信件。但他也要为他的一切行为在基督的审判台前交账，在那里不论大小，都要受审判，不看人的情面。然而，他们怎能信赖他写的信呢？他之所以写那封信，只是因为我们就\Person{保利尼安}及其同伙的非法祝圣一事质问他，正如在那封信的开头他暗示的那样？} 我问，这种盲目是什么意思？他怎么，如俗语所说，沉浸在西梅里安般的黑暗之中？他说我们制造借口，我们没有\Person{厄丕法尼}反对他的信件，他立即补充说：\Quote{他们怎能信赖他写的信？他之所以写那封信，只是因为我们就\Person{保利尼安}及其同伙的非法祝圣一事责问他，正如在那封信的开头他暗示的那样？} 我们没有这样的信！那信又是什么，在开头就提到\Person{保利尼安}？在信的正文中，有你所不敢提及的内容。好啊！你说，你因\Person{保利尼安}的年龄而质问他。但你自己却祝圣一个人为司铎，并派他作为使者和同僚。你竟敢虚伪地称\Person{保利尼安}为少年，然后却派出你自己的少年司铎。同样，你带走了\Place{提里亚}教会的一位执事\Person{特奥塞卡}，立他为司铎，并将武器交在他手中来反对我们，滥用他的口才来伤害我们。只有你才有权践踏教会的权利；无论你做什么，都是教导的标准；而你竟毫不羞愧地挑战\Person{厄丕法尼}，与你一同站在基督的审判台前。这段的下文大意是：他当面责备\Person{厄丕法尼}，说他曾是他的同桌和同屋，并宣称他们从未一起谈论过\Person{奥利金}的观点，而且他以宣誓作证来支持他的话，说：\Quote{他从未显示，天主作证，他曾经怀疑我们的信仰不正确？} 我不愿以刻薄的口吻回答和争辩，以免显得我在指控一位主教作伪证。我们手里有几封\Person{厄丕法尼}的信。一封是写给\Person{若翰}本人的，其他是写给巴勒斯坦的主教们的，还有一封近期写给\Place{罗马}教宗的；在这些信中，他陈述自己曾在众人面前驳斥他的观点，并说他认为对方不值得答复，\Quote{以及整个隐修院，}他说，\Quote{是我们这些卑微者所断言之事的见证。}\par

\SourceRecord{Translated by W.H. Fremantle, G. Lewis and W.G. Martley.；Nicene and Post-Nicene Fathers, Second Series；Vol. 6.；Edited by Philip Schaff and Henry Wace.；Buffalo, NY: Christian Literature Publishing Co.,；1893.；<http://www.newadvent.org/fathers/3004.htm>.}
